\section{Analyses added afterwards}\label{sec:further}

Everything in this section was computed after the results of the second study had been read. It was not fixed
in advance, decides no verdict and carries less weight.

\subsection{The first study, repeated and pooled}\label{sec:replication}

\begin{figure}[!htb]\centering
\begin{tikzpicture}[fig]
\mkdot{cpeople}{1}{7.39}
\node[leg,anchor=base west] at (1.23,7.3) {raters: reliability of the mean of five};
\mkdot{cjev}{6.898}{7.39}
\node[leg,anchor=base west] at (7.128,7.3) {jev-1.14.0};
\mkdia{cdrex}{8.968}{7.39}
\node[leg,anchor=base west] at (9.198,7.3) {drex-max};
\draw[grid] (1.757,0) -- (1.757,6);
\node[tick,anchor=north] at (1.757,-0.1) {0.4};
\draw[grid] (3.471,0) -- (3.471,6);
\node[tick,anchor=north] at (3.471,-0.1) {0.6};
\draw[grid] (5.186,0) -- (5.186,6);
\node[tick,anchor=north] at (5.186,-0.1) {0.8};
\draw[grid] (6.9,0) -- (6.9,6);
\node[tick,anchor=north] at (6.9,-0.1) {1.0};
\draw[grid] (0.9,0.857) -- (6.9,0.857);
\node[tick,anchor=east] at (0.8,0.857) {0.4};
\draw[grid] (0.9,2.571) -- (6.9,2.571);
\node[tick,anchor=east] at (0.8,2.571) {0.6};
\draw[grid] (0.9,4.286) -- (6.9,4.286);
\node[tick,anchor=east] at (0.8,4.286) {0.8};
\draw[grid] (0.9,6) -- (6.9,6);
\node[tick,anchor=east] at (0.8,6) {1.0};
\draw[axis] (0.9,0) -- (6.9,0);
\draw[axis] (0.9,0) -- (0.9,6);
\node[mut,anchor=north] at (3.9,-0.55) {First study (rebuilt for the corpus, levels left in)};
\node[mut,rotate=90,anchor=south] at (-0.08,3) {Second study (plain figures, levels taken out)};
\node[ttl,anchor=base west] at (0.9,6.48) {Do raters agree?};
\node[mut,anchor=base west] at (0.9,6.13) {$r$ across the twelve questions: 0.88};
\draw[ref] (0.9,0) -- (6.9,6);
\node[mut,rotate=45,anchor=south east] at (6.741,5.901) {the same in both};
\draw[grid] (3.471,0) -- (3.471,6);
\draw[grid] (0.9,2.571) -- (6.9,2.571);
\mkdot{cpeople}{2.22}{1.979}
\mkdot{cpeople}{3.183}{2.708}
\mkdot{cpeople}{5.582}{4.234}
\mkdot{cpeople}{5.5}{3.679}
\mkdot{cpeople}{4.8}{4.147}
\mkdot{cpeople}{5.115}{3.742}
\mkdot{cpeople}{2.58}{2.81}
\mkdot{cpeople}{3.075}{2.313}
\mkdot{cpeople}{2.832}{2.346}
\mkdot{cpeople}{3.585}{3.573}
\mkdot{cpeople}{2.817}{2.664}
\mkdot{cpeople}{2.174}{1.004}
\node[leg,anchor=base west] at (2.314,0.924) {Hedging};
\node[leg,anchor=base west] at (3.215,2.233) {Narrative};
\node[leg,anchor=base east] at (2.692,2.266) {Energetic};
\node[leg,anchor=base east] at (2.677,2.584) {Humor};
\node[leg,anchor=base east] at (3.445,3.493) {Calm};
\node[leg,anchor=base east] at (4.66,4.067) {Self-disclosure};
\draw[grid] (9.857,0) -- (9.857,6);
\node[tick,anchor=north] at (9.857,-0.1) {0.4};
\draw[grid] (11.571,0) -- (11.571,6);
\node[tick,anchor=north] at (11.571,-0.1) {0.6};
\draw[grid] (13.286,0) -- (13.286,6);
\node[tick,anchor=north] at (13.286,-0.1) {0.8};
\draw[grid] (15,0) -- (15,6);
\node[tick,anchor=north] at (15,-0.1) {1.0};
\draw[grid] (9,0.857) -- (15,0.857);
\node[tick,anchor=east] at (8.9,0.857) {0.4};
\draw[grid] (9,2.571) -- (15,2.571);
\node[tick,anchor=east] at (8.9,2.571) {0.6};
\draw[grid] (9,4.286) -- (15,4.286);
\node[tick,anchor=east] at (8.9,4.286) {0.8};
\draw[grid] (9,6) -- (15,6);
\node[tick,anchor=east] at (8.9,6) {1.0};
\draw[axis] (9,0) -- (15,0);
\draw[axis] (9,0) -- (9,6);
\node[mut,anchor=north] at (12,-0.55) {First study (rebuilt for the corpus, levels left in)};
\node[mut,rotate=90,anchor=south] at (8.02,3) {Second study (plain figures, levels taken out)};
\node[ttl,anchor=base west] at (9,6.48) {Do the models follow them?};
\node[mut,anchor=base west] at (9,6.13) {$r$ across the twelve: jev-1.14.0 0.95; drex-max 0.90};
\draw[ref] (9,0) -- (15,6);
\node[mut,rotate=45,anchor=south east] at (14.841,5.901) {the same in both};
\draw[grid] (10.714,0) -- (10.714,6);
\draw[grid] (9,1.714) -- (15,1.714);
\mkdot{cjev}{11.533}{2.802}
\mkdot{cjev}{11.98}{2.981}
\mkdot{cjev}{14.07}{4.679}
\mkdot{cjev}{12.896}{3.695}
\mkdot{cjev}{13.219}{4.167}
\mkdot{cjev}{12.926}{4.148}
\mkdot{cjev}{10.762}{2.381}
\mkdot{cjev}{10.428}{1.799}
\mkdot{cjev}{12.541}{3.342}
\mkdot{cjev}{12.09}{3.476}
\mkdot{cjev}{12.409}{2.969}
\mkdot{cjev}{11.535}{2.149}
\mkdia{cdrex}{11.539}{2.445}
\mkdia{cdrex}{11.572}{2.286}
\mkdia{cdrex}{13.947}{4.069}
\mkdia{cdrex}{12.561}{3.458}
\mkdia{cdrex}{13.673}{4.894}
\mkdia{cdrex}{12.911}{3.696}
\mkdia{cdrex}{10.938}{2.731}
\mkdia{cdrex}{11.317}{2.292}
\mkdia{cdrex}{11.995}{2.498}
\mkdia{cdrex}{12.259}{3.354}
\mkdia{cdrex}{11.837}{2.741}
\mkdia{cdrex}{11.49}{2.266}
\end{tikzpicture}%

\caption{The twelve questions rated in both studies: the first study's figure against the second's. Left: the
reliability of the mean of five raters. Right: each model's correlation with that mean. Grey lines: the
thresholds of the rule.}\label{fig:s2replication}
\end{figure}

Twelve questions were rated in both studies, by different people, on different excerpts and in a different
design (Figure~\ref{fig:s2replication}; Table~\ref{tab:s2replication}).
\begin{itemize}
\item \textbf{How well the models follow people repeats closely.} Across the twelve questions the first study's
  correlations and the second's correlate at 0.95 for \jev{} and 0.90 for \drex{}, and no question moves by more
  than 0.07 and 0.10.
\item \textbf{How well raters agree repeats in order, less in size.} The reliabilities correlate at 0.88 across
  the questions. Single questions move by up to 0.13.
\end{itemize}

\begin{table}[!htb]\centering
\caption{The twelve questions rated in both studies, pooled by inverse variance. Both studies' figures are for
the mean of five raters with rater levels left in, so that they estimate the same thing; 95\% intervals. The
rule is applied to the pooled figures.}\label{tab:s2pooled}
{\scriptsize\setlength{\tabcolsep}{3.5pt}
\begin{tabular}{@{}l cc l l l l@{}}
\toprule
 & \multicolumn{3}{l}{Mean of five raters, levels left in} & & & \\
\cmidrule(r){2-4}
Question & First & Second & Pooled & \texttt{jev-1.14.0}, pooled & drex-max, pooled & By the rule, pooled \\
\midrule
Persuasion & 0.45 & 0.46 & 0.46 {\scriptsize[0.39, 0.53]} & 0.59 {\scriptsize[0.53, 0.65]} & 0.58 {\scriptsize[0.52, 0.64]} & borderline \\
Persistence & 0.57 & 0.56 & 0.56 {\scriptsize[0.50, 0.62]} & 0.63 {\scriptsize[0.57, 0.69]} & 0.58 {\scriptsize[0.51, 0.65]} & borderline \\
Positive tone & 0.85 & 0.80 & 0.83 {\scriptsize[0.80, 0.85]} & 0.88 {\scriptsize[0.86, 0.90]} & 0.87 {\scriptsize[0.84, 0.89]} & \textbf{usable}, settled \\
Disappointment & 0.84 & 0.71 & 0.80 {\scriptsize[0.77, 0.83]} & 0.74 {\scriptsize[0.69, 0.79]} & 0.71 {\scriptsize[0.67, 0.75]} & \textbf{usable}, settled \\
Self-disclosure & 0.76 & 0.74 & 0.75 {\scriptsize[0.71, 0.78]} & 0.79 {\scriptsize[0.76, 0.81]} & 0.84 {\scriptsize[0.82, 0.87]} & \textbf{usable}, settled \\
Self-critical & 0.79 & 0.72 & 0.78 {\scriptsize[0.74, 0.82]} & 0.76 {\scriptsize[0.72, 0.81]} & 0.75 {\scriptsize[0.70, 0.79]} & \textbf{usable}, settled \\
Compassion & 0.50 & 0.56 & 0.52 {\scriptsize[0.43, 0.61]} & 0.52 {\scriptsize[0.43, 0.60]} & 0.54 {\scriptsize[0.47, 0.61]} & borderline \\
Narrative & 0.55 & 0.54 & 0.55 {\scriptsize[0.48, 0.62]} & 0.47 {\scriptsize[0.40, 0.55]} & 0.57 {\scriptsize[0.50, 0.64]} & borderline \\
Energetic & 0.53 & 0.49 & 0.51 {\scriptsize[0.46, 0.57]} & 0.69 {\scriptsize[0.65, 0.73]} & 0.63 {\scriptsize[0.59, 0.67]} & borderline, settled \\
Calm & 0.61 & 0.67 & 0.65 {\scriptsize[0.60, 0.69]} & 0.67 {\scriptsize[0.62, 0.72]} & 0.67 {\scriptsize[0.63, 0.72]} & \textbf{usable}, settled \\
Humor & 0.52 & 0.58 & 0.55 {\scriptsize[0.49, 0.61]} & 0.67 {\scriptsize[0.62, 0.72]} & 0.62 {\scriptsize[0.57, 0.67]} & borderline \\
Hedging & 0.45 & 0.37 & 0.42 {\scriptsize[0.33, 0.50]} & 0.57 {\scriptsize[0.51, 0.63]} & 0.57 {\scriptsize[0.52, 0.63]} & borderline \\
\bottomrule
\end{tabular}}

\end{table}

Together the two studies have 320 excerpts for each of these questions (Table~\ref{tab:s2pooled}). Five
questions are usable and settled: the four of the first study and calm. Energetic is settled as borderline: it
is not close to usable for five raters, although both models follow it. The other six still have a line inside
their interval. One question differs between the studies beyond chance: disappointment, the one question that
changed form, was rated more reliably in the first study.

\subsection{How firm the settled verdicts are}\label{sec:firm}

\begin{table}[!htb]\centering
\caption{Standard errors of the primary figures from the choice of recordings (resampling) and of raters (each
rater left out in turn), the 95\% interval from recordings alone and from both, and whether the verdict is
settled with both. The two sources overlap in part, so the interval with both is on the wide
side.}\label{tab:s2raters}
{\scriptsize\setlength{\tabcolsep}{3.2pt}
\begin{tabular}{@{}l l rr l l rr l l@{}}
\toprule
 & & \multicolumn{4}{l}{Reliability of the composite} & \multicolumn{3}{l}{Better model's correlation} & \\
\cmidrule(r){3-6}\cmidrule(r){7-9}
 & Verdict & \multicolumn{2}{l}{SE: rec., raters} & Recordings & Both & \multicolumn{2}{l}{SE: rec., raters} & Both & With both \\
\midrule
F1 & \textbf{usable} & 0.025 & 0.034 & {\scriptsize[0.65, 0.75]} & {\scriptsize[0.61, 0.77]} & 0.031 & 0.032 & {\scriptsize[0.62, 0.81]} & settled \\
F2 & \textbf{usable} & 0.015 & 0.014 & {\scriptsize[0.82, 0.88]} & {\scriptsize[0.81, 0.89]} & 0.028 & 0.021 & {\scriptsize[0.71, 0.85]} & settled \\
F3 & \textbf{usable} & 0.018 & 0.018 & {\scriptsize[0.80, 0.86]} & {\scriptsize[0.78, 0.87]} & 0.025 & 0.023 & {\scriptsize[0.73, 0.86]} & settled \\
F4 & \textbf{usable} & 0.059 & 0.044 & {\scriptsize[0.53, 0.75]} & {\scriptsize[0.50, 0.77]} & 0.093 & 0.052 & {\scriptsize[0.28, 0.69]} & at the line \\
F5 & borderline & 0.055 & 0.057 & {\scriptsize[0.45, 0.68]} & {\scriptsize[0.39, 0.72]} & 0.046 & 0.036 & {\scriptsize[0.55, 0.80]} & at the line \\
F6 & \textbf{usable} & 0.023 & 0.025 & {\scriptsize[0.71, 0.80]} & {\scriptsize[0.69, 0.82]} & 0.024 & 0.021 & {\scriptsize[0.75, 0.87]} & settled \\
F7 & borderline & 0.060 & 0.051 & {\scriptsize[0.44, 0.67]} & {\scriptsize[0.39, 0.70]} & 0.053 & 0.051 & {\scriptsize[0.40, 0.70]} & at the line \\
F8 & unclear rubric & 0.078 & 0.105 & {\scriptsize[0.16, 0.46]} & {\scriptsize[0.04, 0.54]} & 0.052 & 0.067 & {\scriptsize[0.41, 0.75]} & at the line \\
\bottomrule
\end{tabular}}

\end{table}

\paragraph{Raters as well as recordings.} The intervals of Section~\ref{sec:secondresults} describe which
recordings were drawn. Leaving each of the 211 raters out in turn shows that the choice of raters matters about
as much (Table~\ref{tab:s2raters}). With both sources counted the four settled verdicts stay settled; F1 has
the least room, its interval ending at 0.61. Without any one rater F1, F2, F3 and F6 are usable in all 211 runs,
while F4 is usable in 168 and borderline in 43: that is what ``at the line'' means in terms of people.

\begin{table}[!htb]\centering
\caption{Checks of the composite's reliability. Levels taken out: as analysed, and with every rating replaced by
the normal score of its level. Levels left in: from the five rounds of a form taken as five separate panels of
one rater per excerpt, which needs no statistical model; from the variances of the model; and the first study's
figure for the four dimensions that had two markers there. Last column: the verdict when the whole analysis is
run on normal scores.}\label{tab:s2checks}
{\scriptsize\setlength{\tabcolsep}{3.4pt}
\begin{tabular}{@{}l cc l cc l@{}}
\toprule
 & \multicolumn{2}{l}{Levels taken out} & \multicolumn{3}{l}{Levels left in} & \\
\cmidrule(r){2-3}\cmidrule(r){4-6}
Dimension & As analysed & Normal scores & Rounds, no model & Model & First study & On normal scores \\
\midrule
F1 Direct, organized argument & 0.71 & 0.70 & 0.65 {\scriptsize[0.56, 0.72]} & 0.65 & 0.65 & \textbf{usable} \\
F2 Positive tone vs disappointment & 0.86 & 0.86 & 0.86 {\scriptsize[0.81, 0.89]} & 0.85 & 0.88 & \textbf{usable} \\
F3 Self-disclosure, self-criticism & 0.84 & 0.84 & 0.81 {\scriptsize[0.75, 0.85]} & 0.81 & 0.84 & \textbf{usable} \\
F4 Care and regard for others & 0.66 & 0.66 & 0.59 {\scriptsize[0.44, 0.68]} & 0.62 &  & borderline \\
F5 Narrative vs rebuttal & 0.60 & 0.60 & 0.57 {\scriptsize[0.44, 0.68]} & 0.55 &  & borderline \\
F6 Energetic and angry vs calm & 0.77 & 0.77 & 0.72 {\scriptsize[0.65, 0.77]} & 0.72 & 0.70 & \textbf{usable} \\
F7 Humor, imaginative language & 0.59 & 0.61 & 0.49 {\scriptsize[0.32, 0.60]} & 0.53 &  & \textbf{usable} \\
F8 Hedging vs certainty & 0.34 & 0.35 & 0.33 {\scriptsize[0.09, 0.50]} & 0.29 &  & unclear rubric \\
\bottomrule
\end{tabular}}

\end{table}

\paragraph{Without a model, and as ordered categories.} Each round of a form is a complete panel: every excerpt
is read once, by another rater than in the other rounds. The agreement between the five rounds gives the
reliability by a plain analysis of variance, and it matches the model's figure within 0.04 on every dimension
(Table~\ref{tab:s2checks}). When the five levels are treated as ordered categories and not as equally spaced
numbers, no reliability and no correlation moves by more than 0.02, and no settled verdict changes. The two
labels closest to their lines change sides (F4 becomes borderline, F7 usable), which is what ``at the line''
means in terms of method.

\subsection{Where differences between ratings come from}\label{sec:sources}

\begin{figure}[!htb]\centering
\begin{tikzpicture}[fig]
\fill[caccent] (2.9,1.03) rectangle (3.1,1.23);
\node[leg,anchor=base west] at (3.23,1.04) {the excerpt (what raters share)};
\fill[cgone] (7.742,1.03) rectangle (7.942,1.23);
\node[leg,anchor=base west] at (8.072,1.04) {the rater's level};
\fill[cgtwo] (10.736,1.03) rectangle (10.936,1.23);
\node[leg,anchor=base west] at (11.066,1.04) {the rest};
\node[ttl,anchor=base west] at (2.9,0.42) {Sources of the differences between single ratings};
\node[ttl,anchor=base west] at (11.7,0.42) {Mean of five: levels left in};
\node[mut,anchor=base west] at (11.7,0.06) {$\rightarrow$ taken out, as analysed};
\draw[grid] (2.9,0) -- (2.9,-7.36);
\node[tick,anchor=north] at (2.9,-7.48) {0\%};
\draw[grid] (4.975,0) -- (4.975,-7.36);
\node[tick,anchor=north] at (4.975,-7.48) {25\%};
\draw[grid] (7.05,0) -- (7.05,-7.36);
\node[tick,anchor=north] at (7.05,-7.48) {50\%};
\draw[grid] (9.125,0) -- (9.125,-7.36);
\node[tick,anchor=north] at (9.125,-7.48) {75\%};
\draw[grid] (11.2,0) -- (11.2,-7.36);
\node[tick,anchor=north] at (11.2,-7.48) {100\%};
\draw[axis] (2.9,0) -- (2.9,-7.36);
\node[mut,anchor=north] at (7.05,-7.94) {Share of the differences between single ratings of a question};
\node[lbl,anchor=base east] at (2.7,-0.32) {Persuasion};
\fill[caccent] (2.9,-0.36) rectangle (4.081,-0.1);
\node[inbar,text=white,anchor=base west] at (2.98,-0.315) {15\%};
\fill[cgone] (4.116,-0.36) rectangle (6.777,-0.1);
\node[inbar,text=white,anchor=base west] at (4.196,-0.315) {32\%};
\fill[cgtwo] (6.812,-0.36) rectangle (11.2,-0.1);
\node[inbar,text=cink,anchor=base west] at (6.892,-0.315) {53\%};
\node[lbl,anchor=base west] at (11.7,-0.32) {0.46 $\rightarrow$ 0.53};
\node[lbl,anchor=base east] at (2.7,-0.78) {Persistence};
\fill[caccent] (2.9,-0.82) rectangle (4.528,-0.56);
\node[inbar,text=white,anchor=base west] at (2.98,-0.775) {20\%};
\fill[cgone] (4.563,-0.82) rectangle (6.78,-0.56);
\node[inbar,text=white,anchor=base west] at (4.643,-0.775) {27\%};
\fill[cgtwo] (6.815,-0.82) rectangle (11.2,-0.56);
\node[inbar,text=cink,anchor=base west] at (6.895,-0.775) {53\%};
\node[lbl,anchor=base west] at (11.7,-0.78) {0.56 $\rightarrow$ 0.62};
\node[lbl,anchor=base east] at (2.7,-1.24) {Positive tone};
\fill[caccent] (2.9,-1.28) rectangle (6.518,-1.02);
\node[inbar,text=white,anchor=base west] at (2.98,-1.235) {44\%};
\fill[cgone] (6.553,-1.28) rectangle (6.929,-1.02);
\fill[cgtwo] (6.964,-1.28) rectangle (11.2,-1.02);
\node[inbar,text=cink,anchor=base west] at (7.044,-1.235) {51\%};
\node[lbl,anchor=base west] at (11.7,-1.24) {0.80 $\rightarrow$ 0.79};
\node[lbl,anchor=base east] at (2.7,-1.7) {Disappointment};
\fill[caccent] (2.9,-1.74) rectangle (5.589,-1.48);
\node[inbar,text=white,anchor=base west] at (2.98,-1.695) {33\%};
\fill[cgone] (5.624,-1.74) rectangle (6.727,-1.48);
\node[inbar,text=white,anchor=base west] at (5.704,-1.695) {14\%};
\fill[cgtwo] (6.762,-1.74) rectangle (11.2,-1.48);
\node[inbar,text=cink,anchor=base west] at (6.842,-1.695) {53\%};
\node[lbl,anchor=base west] at (11.7,-1.7) {0.71 $\rightarrow$ 0.73};
\node[lbl,anchor=base east] at (2.7,-2.16) {Self-disclosure};
\fill[caccent] (2.9,-2.2) rectangle (5.834,-1.94);
\node[inbar,text=white,anchor=base west] at (2.98,-2.155) {36\%};
\fill[cgone] (5.869,-2.2) rectangle (7.672,-1.94);
\node[inbar,text=white,anchor=base west] at (5.949,-2.155) {22\%};
\fill[cgtwo] (7.707,-2.2) rectangle (11.2,-1.94);
\node[inbar,text=cink,anchor=base west] at (7.787,-2.155) {42\%};
\node[lbl,anchor=base west] at (11.7,-2.16) {0.74 $\rightarrow$ 0.78};
\node[lbl,anchor=base east] at (2.7,-2.62) {Self-critical};
\fill[caccent] (2.9,-2.66) rectangle (5.7,-2.4);
\node[inbar,text=white,anchor=base west] at (2.98,-2.615) {34\%};
\fill[cgone] (5.735,-2.66) rectangle (6.733,-2.4);
\node[inbar,text=white,anchor=base west] at (5.815,-2.615) {12\%};
\fill[cgtwo] (6.768,-2.66) rectangle (11.2,-2.4);
\node[inbar,text=cink,anchor=base west] at (6.848,-2.615) {53\%};
\node[lbl,anchor=base west] at (11.7,-2.62) {0.72 $\rightarrow$ 0.74};
\node[lbl,anchor=base east] at (2.7,-3.08) {Compassion};
\fill[caccent] (2.9,-3.12) rectangle (4.552,-2.86);
\node[inbar,text=white,anchor=base west] at (2.98,-3.075) {20\%};
\fill[cgone] (4.587,-3.12) rectangle (6.937,-2.86);
\node[inbar,text=white,anchor=base west] at (4.667,-3.075) {29\%};
\fill[cgtwo] (6.972,-3.12) rectangle (11.2,-2.86);
\node[inbar,text=cink,anchor=base west] at (7.052,-3.075) {51\%};
\node[lbl,anchor=base west] at (11.7,-3.08) {0.56 $\rightarrow$ 0.63};
\node[lbl,anchor=base east] at (2.7,-3.54) {Respect\textsuperscript{n}};
\fill[caccent] (2.9,-3.58) rectangle (3.811,-3.32);
\node[inbar,text=white,anchor=base west] at (2.98,-3.535) {11\%};
\fill[cgone] (3.846,-3.58) rectangle (5.706,-3.32);
\node[inbar,text=white,anchor=base west] at (3.926,-3.535) {23\%};
\fill[cgtwo] (5.741,-3.58) rectangle (11.2,-3.32);
\node[inbar,text=cink,anchor=base west] at (5.821,-3.535) {66\%};
\node[lbl,anchor=base west] at (11.7,-3.54) {0.39 $\rightarrow$ 0.42};
\node[lbl,anchor=base east] at (2.7,-4) {Narrative};
\fill[caccent] (2.9,-4.04) rectangle (4.434,-3.78);
\node[inbar,text=white,anchor=base west] at (2.98,-3.995) {19\%};
\fill[cgone] (4.469,-4.04) rectangle (5.973,-3.78);
\node[inbar,text=white,anchor=base west] at (4.549,-3.995) {19\%};
\fill[cgtwo] (6.008,-4.04) rectangle (11.2,-3.78);
\node[inbar,text=cink,anchor=base west] at (6.088,-3.995) {63\%};
\node[lbl,anchor=base west] at (11.7,-4) {0.54 $\rightarrow$ 0.57};
\node[lbl,anchor=base east] at (2.7,-4.46) {Rebuttal\textsuperscript{n}};
\fill[caccent] (2.9,-4.5) rectangle (4.042,-4.24);
\node[inbar,text=white,anchor=base west] at (2.98,-4.455) {14\%};
\fill[cgone] (4.077,-4.5) rectangle (6.782,-4.24);
\node[inbar,text=white,anchor=base west] at (4.157,-4.455) {33\%};
\fill[cgtwo] (6.817,-4.5) rectangle (11.2,-4.24);
\node[inbar,text=cink,anchor=base west] at (6.897,-4.455) {53\%};
\node[lbl,anchor=base west] at (11.7,-4.46) {0.45 $\rightarrow$ 0.52};
\node[lbl,anchor=base east] at (2.7,-4.92) {Energetic};
\fill[caccent] (2.9,-4.96) rectangle (4.202,-4.7);
\node[inbar,text=white,anchor=base west] at (2.98,-4.915) {16\%};
\fill[cgone] (4.237,-4.96) rectangle (7.01,-4.7);
\node[inbar,text=white,anchor=base west] at (4.317,-4.915) {34\%};
\fill[cgtwo] (7.045,-4.96) rectangle (11.2,-4.7);
\node[inbar,text=cink,anchor=base west] at (7.125,-4.915) {50\%};
\node[lbl,anchor=base west] at (11.7,-4.92) {0.49 $\rightarrow$ 0.57};
\node[lbl,anchor=base east] at (2.7,-5.38) {Calm};
\fill[caccent] (2.9,-5.42) rectangle (5.296,-5.16);
\node[inbar,text=white,anchor=base west] at (2.98,-5.375) {29\%};
\fill[cgone] (5.331,-5.42) rectangle (7.054,-5.16);
\node[inbar,text=white,anchor=base west] at (5.411,-5.375) {21\%};
\fill[cgtwo] (7.089,-5.42) rectangle (11.2,-5.16);
\node[inbar,text=cink,anchor=base west] at (7.169,-5.375) {50\%};
\node[lbl,anchor=base west] at (11.7,-5.38) {0.67 $\rightarrow$ 0.72};
\node[lbl,anchor=base east] at (2.7,-5.84) {Humor};
\fill[caccent] (2.9,-5.88) rectangle (4.636,-5.62);
\node[inbar,text=white,anchor=base west] at (2.98,-5.835) {21\%};
\fill[cgone] (4.671,-5.88) rectangle (6.346,-5.62);
\node[inbar,text=white,anchor=base west] at (4.751,-5.835) {21\%};
\fill[cgtwo] (6.381,-5.88) rectangle (11.2,-5.62);
\node[inbar,text=cink,anchor=base west] at (6.461,-5.835) {58\%};
\node[lbl,anchor=base west] at (11.7,-5.84) {0.58 $\rightarrow$ 0.61};
\node[lbl,anchor=base east] at (2.7,-6.3) {Imaginative\textsuperscript{n}};
\fill[caccent] (2.9,-6.34) rectangle (3.471,-6.08);
\fill[cgone] (3.506,-6.34) rectangle (6.315,-6.08);
\node[inbar,text=white,anchor=base west] at (3.586,-6.295) {34\%};
\fill[cgtwo] (6.35,-6.34) rectangle (11.2,-6.08);
\node[inbar,text=cink,anchor=base west] at (6.43,-6.295) {58\%};
\node[lbl,anchor=base west] at (11.7,-6.3) {0.28 $\rightarrow$ 0.35};
\node[lbl,anchor=base east] at (2.7,-6.76) {Hedging};
\fill[caccent] (2.9,-6.8) rectangle (3.744,-6.54);
\node[inbar,text=white,anchor=base west] at (2.98,-6.755) {11\%};
\fill[cgone] (3.779,-6.8) rectangle (5.88,-6.54);
\node[inbar,text=white,anchor=base west] at (3.859,-6.755) {26\%};
\fill[cgtwo] (5.915,-6.8) rectangle (11.2,-6.54);
\node[inbar,text=cink,anchor=base west] at (5.995,-6.755) {64\%};
\node[lbl,anchor=base west] at (11.7,-6.76) {0.37 $\rightarrow$ 0.42};
\node[lbl,anchor=base east] at (2.7,-7.22) {Knowledge limits\textsuperscript{n}};
\fill[caccent] (2.9,-7.26) rectangle (3.019,-7);
\fill[cgone] (3.054,-7.26) rectangle (6.06,-7);
\node[inbar,text=white,anchor=base west] at (3.134,-7.215) {37\%};
\fill[cgtwo] (6.095,-7.26) rectangle (11.2,-7);
\node[inbar,text=cink,anchor=base west] at (6.175,-7.215) {62\%};
\node[lbl,anchor=base west] at (11.7,-7.22) {0.09 $\rightarrow$ 0.11};
\end{tikzpicture}%

\caption{Per question, the shares of the differences between two single ratings that come from the excerpt, from
the rater's general level and from the rest. Right: the reliability of the mean of five raters with the levels
left in, as in a design with closed sets, and taken out, as analysed. \textsuperscript{n}: not rated in the
first study.}\label{fig:s2sources}
\end{figure}

A rater's general level, how high or low a person rates a question whatever the excerpt, accounts for 5\%
(positive tone) to 37\% (knowledge limits) of the differences between single ratings, and the excerpt for 44\%
down to 2\% (Figure~\ref{fig:s2sources}). In the first study the level could not be separated from the excerpts
of a set. Taking it out raises the reliability of a composite by 0.01 to 0.06.

\subsection{What the models follow}\label{sec:follow}

\begin{figure}[!htb]\centering
\begin{tikzpicture}[fig]
\mkdot{cjev}{4.7}{1.13}
\node[leg,anchor=base west] at (4.93,1.04) {jev-1.14.0 (upper mark)};
\mkdia{cdrex}{8.486}{1.13}
\node[leg,anchor=base west] at (8.716,1.04) {drex-max (lower mark)};
\mkiv{cjev}{11.858}{12.158}{1.13}
\node[leg,anchor=base west] at (12.238,1.04) {95\% interval};
\draw[grid] (0,0) -- (14.8,0);
\draw[grid] (0,-0.5) -- (14.8,-0.5);
\draw[grid] (0,-1) -- (14.8,-1);
\draw[grid] (0,-1.5) -- (14.8,-1.5);
\draw[grid] (0,-2) -- (14.8,-2);
\draw[grid] (0,-2.5) -- (14.8,-2.5);
\draw[grid] (0,-3) -- (14.8,-3);
\draw[grid] (0,-3.5) -- (14.8,-3.5);
\draw[grid] (0,-4) -- (14.8,-4);
\node[ttl,anchor=base west] at (4.6,0.42) {Within a recording};
\draw[grid] (4.6,0) -- (4.6,-4);
\node[tick,anchor=north] at (4.6,-4.12) {0};
\draw[grid] (5.52,0) -- (5.52,-4);
\node[tick,anchor=north] at (5.52,-4.12) {0.2};
\draw[grid] (6.44,0) -- (6.44,-4);
\node[tick,anchor=north] at (6.44,-4.12) {0.4};
\draw[grid] (7.36,0) -- (7.36,-4);
\node[tick,anchor=north] at (7.36,-4.12) {0.6};
\draw[grid] (8.28,0) -- (8.28,-4);
\node[tick,anchor=north] at (8.28,-4.12) {0.8};
\draw[grid] (9.2,0) -- (9.2,-4);
\node[tick,anchor=north] at (9.2,-4.12) {1.0};
\draw[axis] (4.6,-4) -- (9.2,-4);
\node[mut,anchor=north] at (6.9,-4.58) {Difference between the two passages};
\node[ttl,anchor=base west] at (10.1,0.42) {Between recordings};
\draw[grid] (10.1,0) -- (10.1,-4);
\node[tick,anchor=north] at (10.1,-4.12) {0};
\draw[grid] (11.02,0) -- (11.02,-4);
\node[tick,anchor=north] at (11.02,-4.12) {0.2};
\draw[grid] (11.94,0) -- (11.94,-4);
\node[tick,anchor=north] at (11.94,-4.12) {0.4};
\draw[grid] (12.86,0) -- (12.86,-4);
\node[tick,anchor=north] at (12.86,-4.12) {0.6};
\draw[grid] (13.78,0) -- (13.78,-4);
\node[tick,anchor=north] at (13.78,-4.12) {0.8};
\draw[grid] (14.7,0) -- (14.7,-4);
\node[tick,anchor=north] at (14.7,-4.12) {1.0};
\draw[axis] (10.1,-4) -- (14.7,-4);
\node[mut,anchor=north] at (12.4,-4.58) {Mean of the two passages};
\node[lbl,anchor=base east] at (4.4,-0.34) {F1 Direct, organized argument};
\mkiv{cjev}{6.408}{8.313}{-0.145}
\mkdot{cjev}{7.405}{-0.145}
\mkiv{cdrex}{6.803}{8.46}{-0.355}
\mkdia{cdrex}{7.621}{-0.355}
\mkiv{cjev}{13.751}{14.529}{-0.145}
\mkdot{cjev}{14.145}{-0.145}
\mkiv{cdrex}{13.992}{14.652}{-0.355}
\mkdia{cdrex}{14.317}{-0.355}
\node[lbl,anchor=base east] at (4.4,-0.84) {F2 Positive tone vs disappointment};
\mkiv{cjev}{7.989}{8.949}{-0.645}
\mkdot{cjev}{8.523}{-0.645}
\mkiv{cdrex}{7.337}{8.69}{-0.855}
\mkdia{cdrex}{8.146}{-0.855}
\mkiv{cjev}{13.709}{14.265}{-0.645}
\mkdot{cjev}{14.016}{-0.645}
\mkiv{cdrex}{13.59}{14.209}{-0.855}
\mkdia{cdrex}{13.912}{-0.855}
\node[lbl,anchor=base east] at (4.4,-1.34) {F3 Self-disclosure, self-criticism};
\mkiv{cjev}{7.404}{8.789}{-1.145}
\mkdot{cjev}{8.087}{-1.145}
\mkiv{cdrex}{6.782}{8.456}{-1.355}
\mkdia{cdrex}{7.64}{-1.355}
\mkiv{cjev}{14.014}{14.534}{-1.145}
\mkdot{cjev}{14.309}{-1.145}
\mkiv{cdrex}{14.262}{14.575}{-1.355}
\mkdia{cdrex}{14.426}{-1.355}
\node[lbl,anchor=base east] at (4.4,-1.84) {F4 Care and regard for others};
\mkiv{cjev}{6.637}{9.027}{-1.645}
\mkdot{cjev}{7.973}{-1.645}
\mkiv{cdrex}{6.438}{9.018}{-1.855}
\mkdia{cdrex}{7.93}{-1.855}
\mkiv{cjev}{11.495}{13.694}{-1.645}
\mkdot{cjev}{12.745}{-1.645}
\mkiv{cdrex}{11.593}{13.626}{-1.855}
\mkdia{cdrex}{12.751}{-1.855}
\node[lbl,anchor=base east] at (4.4,-2.34) {F5 Narrative vs rebuttal};
\mkiv{cjev}{7.66}{9.2}{-2.145}
\mkdot{cjev}{8.879}{-2.145}
\mkiv{cdrex}{7.444}{9.2}{-2.355}
\mkdia{cdrex}{8.725}{-2.355}
\mkiv{cjev}{13.951}{14.662}{-2.145}
\mkdot{cjev}{14.345}{-2.145}
\mkiv{cdrex}{13.488}{14.441}{-2.355}
\mkdia{cdrex}{14.026}{-2.355}
\node[lbl,anchor=base east] at (4.4,-2.84) {F6 Energetic and angry vs calm};
\mkiv{cjev}{7.633}{9.2}{-2.645}
\mkdot{cjev}{8.626}{-2.645}
\mkiv{cdrex}{7.282}{9.2}{-2.855}
\mkdia{cdrex}{8.244}{-2.855}
\mkiv{cjev}{14.37}{14.635}{-2.645}
\mkdot{cjev}{14.506}{-2.645}
\mkiv{cdrex}{13.895}{14.43}{-2.855}
\mkdia{cdrex}{14.181}{-2.855}
\node[lbl,anchor=base east] at (4.4,-3.34) {F7 Humor, imaginative language};
\mkiv{cjev}{6.063}{8.695}{-3.145}
\mkdot{cjev}{7.351}{-3.145}
\mkiv{cdrex}{4.672}{7.556}{-3.355}
\mkdia{cdrex}{6.153}{-3.355}
\mkiv{cjev}{13.26}{14.419}{-3.145}
\mkdot{cjev}{13.825}{-3.145}
\mkiv{cdrex}{12.949}{14.278}{-3.355}
\mkdia{cdrex}{13.675}{-3.355}
\node[lbl,anchor=base east] at (4.4,-3.84) {F8 Hedging vs certainty};
\mkiv{cjev}{8.221}{9.2}{-3.645}
\mkdot{cjev}{9.2}{-3.645}
\mkiv{cdrex}{8.424}{9.2}{-3.855}
\mkdia{cdrex}{9.2}{-3.855}
\mkiv{cjev}{12.626}{14.7}{-3.645}
\mkdot{cjev}{13.639}{-3.645}
\mkiv{cdrex}{13.258}{14.7}{-3.855}
\mkdia{cdrex}{14.157}{-3.855}
\end{tikzpicture}%

\caption{Do the models follow the difference between two passages of the same recording? Left: the correlation
of the difference in a model's dimension score with the difference in the panel's consensus, between the two
passages of a recording. Right: the same for the mean of the two passages, between recordings. Rater noise
divided out; 95\% intervals from resampling recordings.}\label{fig:s2within}
\end{figure}

\paragraph{The passage, not only the recording.} A model could follow people by knowing only how a recording
sounds in general. The two passages of each recording test this (Figure~\ref{fig:s2within};
Table~\ref{tab:s2within}). As observed, \jev{}'s within-recording correlations are 0.41 to 0.74, each with an
interval clear of zero. With rater noise divided out, within is close to between: 0.79 against 0.84 on average
for \jev{}, 0.73 against 0.83 for \drex{}. On F1 and F3 part of what the model follows belongs to the
recording.

\begin{figure}[!htb]\centering
\begin{tikzpicture}[fig]
\node[ttl,anchor=base west] at (1,0.95) {jev-1.14.0};
\node[lbl,anchor=base east] at (0.88,-0.52) {F1};
\node[lbl,anchor=base] at (1.43,0.2) {F1};
\fill[cjev!59] (1.02,-0.02) rectangle (1.84,-0.84);
\draw[cink,line width=0.7pt] (1.02,-0.02) rectangle (1.84,-0.84);
\node[inner sep=0pt,font=\sffamily\scriptsize,text=cink] at (1.43,-0.43) {0.69};
\fill[cjev!12] (1.88,-0.02) rectangle (2.7,-0.84);
\node[inner sep=0pt,font=\sffamily\scriptsize,text=cink] at (2.29,-0.43) {0.15};
\fill[cgone!11] (2.74,-0.02) rectangle (3.56,-0.84);
\node[inner sep=0pt,font=\sffamily\scriptsize,text=cink] at (3.15,-0.43) {\textminus 0.13};
\fill[cjev!13] (3.6,-0.02) rectangle (4.42,-0.84);
\node[inner sep=0pt,font=\sffamily\scriptsize,text=cink] at (4.01,-0.43) {0.15};
\fill[cjev!1] (4.46,-0.02) rectangle (5.28,-0.84);
\node[inner sep=0pt,font=\sffamily\scriptsize,text=cink] at (4.87,-0.43) {0.01};
\fill[cgone!41] (5.32,-0.02) rectangle (6.14,-0.84);
\node[inner sep=0pt,font=\sffamily\scriptsize,text=cink] at (5.73,-0.43) {\textminus 0.48};
\fill[cgone!7] (6.18,-0.02) rectangle (7,-0.84);
\node[inner sep=0pt,font=\sffamily\scriptsize,text=cink] at (6.59,-0.43) {\textminus 0.08};
\fill[cgone!10] (7.04,-0.02) rectangle (7.86,-0.84);
\node[inner sep=0pt,font=\sffamily\scriptsize,text=cink] at (7.45,-0.43) {\textminus 0.12};
\node[lbl,anchor=base east] at (0.88,-1.38) {F2};
\node[lbl,anchor=base] at (2.29,0.2) {F2};
\fill[cgone!12] (1.02,-0.88) rectangle (1.84,-1.7);
\node[inner sep=0pt,font=\sffamily\scriptsize,text=cink] at (1.43,-1.29) {\textminus 0.15};
\fill[cjev!67] (1.88,-0.88) rectangle (2.7,-1.7);
\draw[cink,line width=0.7pt] (1.88,-0.88) rectangle (2.7,-1.7);
\node[inner sep=0pt,font=\sffamily\scriptsize,text=cink] at (2.29,-1.29) {0.79};
\fill[cgone!35] (2.74,-0.88) rectangle (3.56,-1.7);
\node[inner sep=0pt,font=\sffamily\scriptsize,text=cink] at (3.15,-1.29) {\textminus 0.42};
\fill[cgone!7] (3.6,-0.88) rectangle (4.42,-1.7);
\node[inner sep=0pt,font=\sffamily\scriptsize,text=cink] at (4.01,-1.29) {\textminus 0.08};
\fill[cjev!2] (4.46,-0.88) rectangle (5.28,-1.7);
\node[inner sep=0pt,font=\sffamily\scriptsize,text=cink] at (4.87,-1.29) {0.02};
\fill[cgone!5] (5.32,-0.88) rectangle (6.14,-1.7);
\node[inner sep=0pt,font=\sffamily\scriptsize,text=cink] at (5.73,-1.29) {\textminus 0.06};
\fill[cjev!1] (6.18,-0.88) rectangle (7,-1.7);
\node[inner sep=0pt,font=\sffamily\scriptsize,text=cink] at (6.59,-1.29) {0.01};
\fill[cgone!7] (7.04,-0.88) rectangle (7.86,-1.7);
\node[inner sep=0pt,font=\sffamily\scriptsize,text=cink] at (7.45,-1.29) {\textminus 0.09};
\node[lbl,anchor=base east] at (0.88,-2.24) {F3};
\node[lbl,anchor=base] at (3.15,0.2) {F3};
\fill[cgone!13] (1.02,-1.74) rectangle (1.84,-2.56);
\node[inner sep=0pt,font=\sffamily\scriptsize,text=cink] at (1.43,-2.15) {\textminus 0.16};
\fill[cgone!44] (1.88,-1.74) rectangle (2.7,-2.56);
\node[inner sep=0pt,font=\sffamily\scriptsize,text=cink] at (2.29,-2.15) {\textminus 0.52};
\fill[cjev!68] (2.74,-1.74) rectangle (3.56,-2.56);
\draw[cink,line width=0.7pt] (2.74,-1.74) rectangle (3.56,-2.56);
\node[inner sep=0pt,font=\sffamily\scriptsize,text=cink] at (3.15,-2.15) {0.80};
\fill[cgone!10] (3.6,-1.74) rectangle (4.42,-2.56);
\node[inner sep=0pt,font=\sffamily\scriptsize,text=cink] at (4.01,-2.15) {\textminus 0.12};
\fill[cgone!4] (4.46,-1.74) rectangle (5.28,-2.56);
\node[inner sep=0pt,font=\sffamily\scriptsize,text=cink] at (4.87,-2.15) {\textminus 0.05};
\fill[cjev!18] (5.32,-1.74) rectangle (6.14,-2.56);
\node[inner sep=0pt,font=\sffamily\scriptsize,text=cink] at (5.73,-2.15) {0.21};
\fill[cgone!16] (6.18,-1.74) rectangle (7,-2.56);
\node[inner sep=0pt,font=\sffamily\scriptsize,text=cink] at (6.59,-2.15) {\textminus 0.18};
\fill[cjev!3] (7.04,-1.74) rectangle (7.86,-2.56);
\node[inner sep=0pt,font=\sffamily\scriptsize,text=cink] at (7.45,-2.15) {0.03};
\node[lbl,anchor=base east] at (0.88,-3.1) {F4};
\node[lbl,anchor=base] at (4.01,0.2) {F4};
\fill[cgone!4] (1.02,-2.6) rectangle (1.84,-3.42);
\node[inner sep=0pt,font=\sffamily\scriptsize,text=cink] at (1.43,-3.01) {\textminus 0.05};
\fill[cgone!10] (1.88,-2.6) rectangle (2.7,-3.42);
\node[inner sep=0pt,font=\sffamily\scriptsize,text=cink] at (2.29,-3.01) {\textminus 0.12};
\fill[cjev!17] (2.74,-2.6) rectangle (3.56,-3.42);
\node[inner sep=0pt,font=\sffamily\scriptsize,text=cink] at (3.15,-3.01) {0.19};
\fill[cjev!43] (3.6,-2.6) rectangle (4.42,-3.42);
\draw[cink,line width=0.7pt] (3.6,-2.6) rectangle (4.42,-3.42);
\node[inner sep=0pt,font=\sffamily\scriptsize,text=cink] at (4.01,-3.01) {0.51};
\fill[cgone!20] (4.46,-2.6) rectangle (5.28,-3.42);
\node[inner sep=0pt,font=\sffamily\scriptsize,text=cink] at (4.87,-3.01) {\textminus 0.24};
\fill[cjev!4] (5.32,-2.6) rectangle (6.14,-3.42);
\node[inner sep=0pt,font=\sffamily\scriptsize,text=cink] at (5.73,-3.01) {0.05};
\fill[cgone!4] (6.18,-2.6) rectangle (7,-3.42);
\node[inner sep=0pt,font=\sffamily\scriptsize,text=cink] at (6.59,-3.01) {\textminus 0.04};
\fill[cjev!11] (7.04,-2.6) rectangle (7.86,-3.42);
\node[inner sep=0pt,font=\sffamily\scriptsize,text=cink] at (7.45,-3.01) {0.13};
\node[lbl,anchor=base east] at (0.88,-3.96) {F5};
\node[lbl,anchor=base] at (4.87,0.2) {F5};
\fill[cgone!11] (1.02,-3.46) rectangle (1.84,-4.28);
\node[inner sep=0pt,font=\sffamily\scriptsize,text=cink] at (1.43,-3.87) {\textminus 0.13};
\fill[cjev!25] (1.88,-3.46) rectangle (2.7,-4.28);
\node[inner sep=0pt,font=\sffamily\scriptsize,text=cink] at (2.29,-3.87) {0.29};
\fill[cgone!9] (2.74,-3.46) rectangle (3.56,-4.28);
\node[inner sep=0pt,font=\sffamily\scriptsize,text=cink] at (3.15,-3.87) {\textminus 0.11};
\fill[cgone!6] (3.6,-3.46) rectangle (4.42,-4.28);
\node[inner sep=0pt,font=\sffamily\scriptsize,text=cink] at (4.01,-3.87) {\textminus 0.07};
\fill[cjev!60] (4.46,-3.46) rectangle (5.28,-4.28);
\draw[cink,line width=0.7pt] (4.46,-3.46) rectangle (5.28,-4.28);
\node[inner sep=0pt,font=\sffamily\scriptsize,text=cink] at (4.87,-3.87) {0.70};
\fill[cgone!15] (5.32,-3.46) rectangle (6.14,-4.28);
\node[inner sep=0pt,font=\sffamily\scriptsize,text=cink] at (5.73,-3.87) {\textminus 0.17};
\fill[cjev!0] (6.18,-3.46) rectangle (7,-4.28);
\node[inner sep=0pt,font=\sffamily\scriptsize,text=cink] at (6.59,-3.87) {0.00};
\fill[cgone!3] (7.04,-3.46) rectangle (7.86,-4.28);
\node[inner sep=0pt,font=\sffamily\scriptsize,text=cink] at (7.45,-3.87) {\textminus 0.03};
\node[lbl,anchor=base east] at (0.88,-4.82) {F6};
\node[lbl,anchor=base] at (5.73,0.2) {F6};
\fill[cgone!31] (1.02,-4.32) rectangle (1.84,-5.14);
\node[inner sep=0pt,font=\sffamily\scriptsize,text=cink] at (1.43,-4.73) {\textminus 0.36};
\fill[cgone!37] (1.88,-4.32) rectangle (2.7,-5.14);
\node[inner sep=0pt,font=\sffamily\scriptsize,text=cink] at (2.29,-4.73) {\textminus 0.44};
\fill[cjev!9] (2.74,-4.32) rectangle (3.56,-5.14);
\node[inner sep=0pt,font=\sffamily\scriptsize,text=cink] at (3.15,-4.73) {0.11};
\fill[cgone!38] (3.6,-4.32) rectangle (4.42,-5.14);
\node[inner sep=0pt,font=\sffamily\scriptsize,text=cink] at (4.01,-4.73) {\textminus 0.45};
\fill[cgone!25] (4.46,-4.32) rectangle (5.28,-5.14);
\node[inner sep=0pt,font=\sffamily\scriptsize,text=cink] at (4.87,-4.73) {\textminus 0.29};
\fill[cjev!70] (5.32,-4.32) rectangle (6.14,-5.14);
\draw[cink,line width=0.7pt] (5.32,-4.32) rectangle (6.14,-5.14);
\node[inner sep=0pt,font=\sffamily\scriptsize,text=cink] at (5.73,-4.73) {0.82};
\fill[cjev!8] (6.18,-4.32) rectangle (7,-5.14);
\node[inner sep=0pt,font=\sffamily\scriptsize,text=cink] at (6.59,-4.73) {0.09};
\fill[cgone!11] (7.04,-4.32) rectangle (7.86,-5.14);
\node[inner sep=0pt,font=\sffamily\scriptsize,text=cink] at (7.45,-4.73) {\textminus 0.13};
\node[lbl,anchor=base east] at (0.88,-5.68) {F7};
\node[lbl,anchor=base] at (6.59,0.2) {F7};
\fill[cgone!23] (1.02,-5.18) rectangle (1.84,-6);
\node[inner sep=0pt,font=\sffamily\scriptsize,text=cink] at (1.43,-5.59) {\textminus 0.27};
\fill[cjev!17] (1.88,-5.18) rectangle (2.7,-6);
\node[inner sep=0pt,font=\sffamily\scriptsize,text=cink] at (2.29,-5.59) {0.20};
\fill[cgone!16] (2.74,-5.18) rectangle (3.56,-6);
\node[inner sep=0pt,font=\sffamily\scriptsize,text=cink] at (3.15,-5.59) {\textminus 0.19};
\fill[cgone!14] (3.6,-5.18) rectangle (4.42,-6);
\node[inner sep=0pt,font=\sffamily\scriptsize,text=cink] at (4.01,-5.59) {\textminus 0.16};
\fill[cgone!14] (4.46,-5.18) rectangle (5.28,-6);
\node[inner sep=0pt,font=\sffamily\scriptsize,text=cink] at (4.87,-5.59) {\textminus 0.16};
\fill[cjev!30] (5.32,-5.18) rectangle (6.14,-6);
\node[inner sep=0pt,font=\sffamily\scriptsize,text=cink] at (5.73,-5.59) {0.35};
\fill[cjev!48] (6.18,-5.18) rectangle (7,-6);
\draw[cink,line width=0.7pt] (6.18,-5.18) rectangle (7,-6);
\node[inner sep=0pt,font=\sffamily\scriptsize,text=cink] at (6.59,-5.59) {0.57};
\fill[cgone!1] (7.04,-5.18) rectangle (7.86,-6);
\node[inner sep=0pt,font=\sffamily\scriptsize,text=cink] at (7.45,-5.59) {\textminus 0.01};
\node[lbl,anchor=base east] at (0.88,-6.54) {F8};
\node[lbl,anchor=base] at (7.45,0.2) {F8};
\fill[cgone!30] (1.02,-6.04) rectangle (1.84,-6.86);
\node[inner sep=0pt,font=\sffamily\scriptsize,text=cink] at (1.43,-6.45) {\textminus 0.35};
\fill[cjev!12] (1.88,-6.04) rectangle (2.7,-6.86);
\node[inner sep=0pt,font=\sffamily\scriptsize,text=cink] at (2.29,-6.45) {0.14};
\fill[cgone!7] (2.74,-6.04) rectangle (3.56,-6.86);
\node[inner sep=0pt,font=\sffamily\scriptsize,text=cink] at (3.15,-6.45) {\textminus 0.08};
\fill[cjev!13] (3.6,-6.04) rectangle (4.42,-6.86);
\node[inner sep=0pt,font=\sffamily\scriptsize,text=cink] at (4.01,-6.45) {0.15};
\fill[cgone!4] (4.46,-6.04) rectangle (5.28,-6.86);
\node[inner sep=0pt,font=\sffamily\scriptsize,text=cink] at (4.87,-6.45) {\textminus 0.05};
\fill[cjev!5] (5.32,-6.04) rectangle (6.14,-6.86);
\node[inner sep=0pt,font=\sffamily\scriptsize,text=cink] at (5.73,-6.45) {0.06};
\fill[cjev!2] (6.18,-6.04) rectangle (7,-6.86);
\node[inner sep=0pt,font=\sffamily\scriptsize,text=cink] at (6.59,-6.45) {0.02};
\fill[cjev!44] (7.04,-6.04) rectangle (7.86,-6.86);
\draw[cink,line width=0.7pt] (7.04,-6.04) rectangle (7.86,-6.86);
\node[inner sep=0pt,font=\sffamily\scriptsize,text=cink] at (7.45,-6.45) {0.52};
\node[mut,anchor=north] at (4.44,-7.1) {People's composite of the dimension};
\node[ttl,anchor=base west] at (8.9,0.95) {drex-max};
\node[lbl,anchor=base east] at (8.78,-0.52) {F1};
\node[lbl,anchor=base] at (9.33,0.2) {F1};
\fill[cdrex!61] (8.92,-0.02) rectangle (9.74,-0.84);
\draw[cink,line width=0.7pt] (8.92,-0.02) rectangle (9.74,-0.84);
\node[inner sep=0pt,font=\sffamily\scriptsize,text=cink] at (9.33,-0.43) {0.72};
\fill[cdrex!7] (9.78,-0.02) rectangle (10.6,-0.84);
\node[inner sep=0pt,font=\sffamily\scriptsize,text=cink] at (10.19,-0.43) {0.08};
\fill[cgone!2] (10.64,-0.02) rectangle (11.46,-0.84);
\node[inner sep=0pt,font=\sffamily\scriptsize,text=cink] at (11.05,-0.43) {\textminus 0.03};
\fill[cdrex!23] (11.5,-0.02) rectangle (12.32,-0.84);
\node[inner sep=0pt,font=\sffamily\scriptsize,text=cink] at (11.91,-0.43) {0.27};
\fill[cgone!4] (12.36,-0.02) rectangle (13.18,-0.84);
\node[inner sep=0pt,font=\sffamily\scriptsize,text=cink] at (12.77,-0.43) {\textminus 0.05};
\fill[cgone!47] (13.22,-0.02) rectangle (14.04,-0.84);
\node[inner sep=0pt,font=\sffamily\scriptsize,text=cink] at (13.63,-0.43) {\textminus 0.55};
\fill[cgone!9] (14.08,-0.02) rectangle (14.9,-0.84);
\node[inner sep=0pt,font=\sffamily\scriptsize,text=cink] at (14.49,-0.43) {\textminus 0.11};
\fill[cgone!5] (14.94,-0.02) rectangle (15.76,-0.84);
\node[inner sep=0pt,font=\sffamily\scriptsize,text=cink] at (15.35,-0.43) {\textminus 0.06};
\node[lbl,anchor=base east] at (8.78,-1.38) {F2};
\node[lbl,anchor=base] at (10.19,0.2) {F2};
\fill[cgone!13] (8.92,-0.88) rectangle (9.74,-1.7);
\node[inner sep=0pt,font=\sffamily\scriptsize,text=cink] at (9.33,-1.29) {\textminus 0.16};
\fill[cdrex!64] (9.78,-0.88) rectangle (10.6,-1.7);
\draw[cink,line width=0.7pt] (9.78,-0.88) rectangle (10.6,-1.7);
\node[inner sep=0pt,font=\sffamily\scriptsize,text=cink] at (10.19,-1.29) {0.76};
\fill[cgone!24] (10.64,-0.88) rectangle (11.46,-1.7);
\node[inner sep=0pt,font=\sffamily\scriptsize,text=cink] at (11.05,-1.29) {\textminus 0.29};
\fill[cdrex!3] (11.5,-0.88) rectangle (12.32,-1.7);
\node[inner sep=0pt,font=\sffamily\scriptsize,text=cink] at (11.91,-1.29) {0.03};
\fill[cdrex!3] (12.36,-0.88) rectangle (13.18,-1.7);
\node[inner sep=0pt,font=\sffamily\scriptsize,text=cink] at (12.77,-1.29) {0.04};
\fill[cgone!9] (13.22,-0.88) rectangle (14.04,-1.7);
\node[inner sep=0pt,font=\sffamily\scriptsize,text=cink] at (13.63,-1.29) {\textminus 0.10};
\fill[cdrex!1] (14.08,-0.88) rectangle (14.9,-1.7);
\node[inner sep=0pt,font=\sffamily\scriptsize,text=cink] at (14.49,-1.29) {0.01};
\fill[cgone!6] (14.94,-0.88) rectangle (15.76,-1.7);
\node[inner sep=0pt,font=\sffamily\scriptsize,text=cink] at (15.35,-1.29) {\textminus 0.07};
\node[lbl,anchor=base east] at (8.78,-2.24) {F3};
\node[lbl,anchor=base] at (11.05,0.2) {F3};
\fill[cgone!5] (8.92,-1.74) rectangle (9.74,-2.56);
\node[inner sep=0pt,font=\sffamily\scriptsize,text=cink] at (9.33,-2.15) {\textminus 0.06};
\fill[cgone!47] (9.78,-1.74) rectangle (10.6,-2.56);
\node[inner sep=0pt,font=\sffamily\scriptsize,text=cink] at (10.19,-2.15) {\textminus 0.56};
\fill[cdrex!69] (10.64,-1.74) rectangle (11.46,-2.56);
\draw[cink,line width=0.7pt] (10.64,-1.74) rectangle (11.46,-2.56);
\node[inner sep=0pt,font=\sffamily\scriptsize,text=cink] at (11.05,-2.15) {0.81};
\fill[cdrex!1] (11.5,-1.74) rectangle (12.32,-2.56);
\node[inner sep=0pt,font=\sffamily\scriptsize,text=cink] at (11.91,-2.15) {0.01};
\fill[cdrex!2] (12.36,-1.74) rectangle (13.18,-2.56);
\node[inner sep=0pt,font=\sffamily\scriptsize,text=cink] at (12.77,-2.15) {0.03};
\fill[cdrex!10] (13.22,-1.74) rectangle (14.04,-2.56);
\node[inner sep=0pt,font=\sffamily\scriptsize,text=cink] at (13.63,-2.15) {0.12};
\fill[cgone!18] (14.08,-1.74) rectangle (14.9,-2.56);
\node[inner sep=0pt,font=\sffamily\scriptsize,text=cink] at (14.49,-2.15) {\textminus 0.21};
\fill[cdrex!7] (14.94,-1.74) rectangle (15.76,-2.56);
\node[inner sep=0pt,font=\sffamily\scriptsize,text=cink] at (15.35,-2.15) {0.09};
\node[lbl,anchor=base east] at (8.78,-3.1) {F4};
\node[lbl,anchor=base] at (11.91,0.2) {F4};
\fill[cdrex!0] (8.92,-2.6) rectangle (9.74,-3.42);
\node[inner sep=0pt,font=\sffamily\scriptsize,text=cink] at (9.33,-3.01) {0.00};
\fill[cgone!20] (9.78,-2.6) rectangle (10.6,-3.42);
\node[inner sep=0pt,font=\sffamily\scriptsize,text=cink] at (10.19,-3.01) {\textminus 0.23};
\fill[cdrex!26] (10.64,-2.6) rectangle (11.46,-3.42);
\node[inner sep=0pt,font=\sffamily\scriptsize,text=cink] at (11.05,-3.01) {0.30};
\fill[cdrex!43] (11.5,-2.6) rectangle (12.32,-3.42);
\draw[cink,line width=0.7pt] (11.5,-2.6) rectangle (12.32,-3.42);
\node[inner sep=0pt,font=\sffamily\scriptsize,text=cink] at (11.91,-3.01) {0.51};
\fill[cgone!19] (12.36,-2.6) rectangle (13.18,-3.42);
\node[inner sep=0pt,font=\sffamily\scriptsize,text=cink] at (12.77,-3.01) {\textminus 0.22};
\fill[cdrex!3] (13.22,-2.6) rectangle (14.04,-3.42);
\node[inner sep=0pt,font=\sffamily\scriptsize,text=cink] at (13.63,-3.01) {0.04};
\fill[cgone!5] (14.08,-2.6) rectangle (14.9,-3.42);
\node[inner sep=0pt,font=\sffamily\scriptsize,text=cink] at (14.49,-3.01) {\textminus 0.06};
\fill[cdrex!6] (14.94,-2.6) rectangle (15.76,-3.42);
\node[inner sep=0pt,font=\sffamily\scriptsize,text=cink] at (15.35,-3.01) {0.07};
\node[lbl,anchor=base east] at (8.78,-3.96) {F5};
\node[lbl,anchor=base] at (12.77,0.2) {F5};
\fill[cgone!1] (8.92,-3.46) rectangle (9.74,-4.28);
\node[inner sep=0pt,font=\sffamily\scriptsize,text=cink] at (9.33,-3.87) {\textminus 0.01};
\fill[cdrex!21] (9.78,-3.46) rectangle (10.6,-4.28);
\node[inner sep=0pt,font=\sffamily\scriptsize,text=cink] at (10.19,-3.87) {0.25};
\fill[cgone!7] (10.64,-3.46) rectangle (11.46,-4.28);
\node[inner sep=0pt,font=\sffamily\scriptsize,text=cink] at (11.05,-3.87) {\textminus 0.09};
\fill[cdrex!3] (11.5,-3.46) rectangle (12.32,-4.28);
\node[inner sep=0pt,font=\sffamily\scriptsize,text=cink] at (11.91,-3.87) {0.04};
\fill[cdrex!55] (12.36,-3.46) rectangle (13.18,-4.28);
\draw[cink,line width=0.7pt] (12.36,-3.46) rectangle (13.18,-4.28);
\node[inner sep=0pt,font=\sffamily\scriptsize,text=cink] at (12.77,-3.87) {0.65};
\fill[cgone!24] (13.22,-3.46) rectangle (14.04,-4.28);
\node[inner sep=0pt,font=\sffamily\scriptsize,text=cink] at (13.63,-3.87) {\textminus 0.28};
\fill[cdrex!5] (14.08,-3.46) rectangle (14.9,-4.28);
\node[inner sep=0pt,font=\sffamily\scriptsize,text=cink] at (14.49,-3.87) {0.06};
\fill[cdrex!3] (14.94,-3.46) rectangle (15.76,-4.28);
\node[inner sep=0pt,font=\sffamily\scriptsize,text=cink] at (15.35,-3.87) {0.04};
\node[lbl,anchor=base east] at (8.78,-4.82) {F6};
\node[lbl,anchor=base] at (13.63,0.2) {F6};
\fill[cgone!25] (8.92,-4.32) rectangle (9.74,-5.14);
\node[inner sep=0pt,font=\sffamily\scriptsize,text=cink] at (9.33,-4.73) {\textminus 0.29};
\fill[cgone!42] (9.78,-4.32) rectangle (10.6,-5.14);
\node[inner sep=0pt,font=\sffamily\scriptsize,text=cink] at (10.19,-4.73) {\textminus 0.49};
\fill[cdrex!15] (10.64,-4.32) rectangle (11.46,-5.14);
\node[inner sep=0pt,font=\sffamily\scriptsize,text=cink] at (11.05,-4.73) {0.18};
\fill[cgone!32] (11.5,-4.32) rectangle (12.32,-5.14);
\node[inner sep=0pt,font=\sffamily\scriptsize,text=cink] at (11.91,-4.73) {\textminus 0.38};
\fill[cgone!19] (12.36,-4.32) rectangle (13.18,-5.14);
\node[inner sep=0pt,font=\sffamily\scriptsize,text=cink] at (12.77,-4.73) {\textminus 0.22};
\fill[cdrex!64] (13.22,-4.32) rectangle (14.04,-5.14);
\draw[cink,line width=0.7pt] (13.22,-4.32) rectangle (14.04,-5.14);
\node[inner sep=0pt,font=\sffamily\scriptsize,text=cink] at (13.63,-4.73) {0.75};
\fill[cdrex!7] (14.08,-4.32) rectangle (14.9,-5.14);
\node[inner sep=0pt,font=\sffamily\scriptsize,text=cink] at (14.49,-4.73) {0.08};
\fill[cgone!10] (14.94,-4.32) rectangle (15.76,-5.14);
\node[inner sep=0pt,font=\sffamily\scriptsize,text=cink] at (15.35,-4.73) {\textminus 0.12};
\node[lbl,anchor=base east] at (8.78,-5.68) {F7};
\node[lbl,anchor=base] at (14.49,0.2) {F7};
\fill[cgone!9] (8.92,-5.18) rectangle (9.74,-6);
\node[inner sep=0pt,font=\sffamily\scriptsize,text=cink] at (9.33,-5.59) {\textminus 0.10};
\fill[cdrex!11] (9.78,-5.18) rectangle (10.6,-6);
\node[inner sep=0pt,font=\sffamily\scriptsize,text=cink] at (10.19,-5.59) {0.13};
\fill[cgone!8] (10.64,-5.18) rectangle (11.46,-6);
\node[inner sep=0pt,font=\sffamily\scriptsize,text=cink] at (11.05,-5.59) {\textminus 0.09};
\fill[cgone!6] (11.5,-5.18) rectangle (12.32,-6);
\node[inner sep=0pt,font=\sffamily\scriptsize,text=cink] at (11.91,-5.59) {\textminus 0.08};
\fill[cgone!20] (12.36,-5.18) rectangle (13.18,-6);
\node[inner sep=0pt,font=\sffamily\scriptsize,text=cink] at (12.77,-5.59) {\textminus 0.24};
\fill[cdrex!24] (13.22,-5.18) rectangle (14.04,-6);
\node[inner sep=0pt,font=\sffamily\scriptsize,text=cink] at (13.63,-5.59) {0.28};
\fill[cdrex!42] (14.08,-5.18) rectangle (14.9,-6);
\draw[cink,line width=0.7pt] (14.08,-5.18) rectangle (14.9,-6);
\node[inner sep=0pt,font=\sffamily\scriptsize,text=cink] at (14.49,-5.59) {0.50};
\fill[cgone!6] (14.94,-5.18) rectangle (15.76,-6);
\node[inner sep=0pt,font=\sffamily\scriptsize,text=cink] at (15.35,-5.59) {\textminus 0.07};
\node[lbl,anchor=base east] at (8.78,-6.54) {F8};
\node[lbl,anchor=base] at (15.35,0.2) {F8};
\fill[cgone!21] (8.92,-6.04) rectangle (9.74,-6.86);
\node[inner sep=0pt,font=\sffamily\scriptsize,text=cink] at (9.33,-6.45) {\textminus 0.25};
\fill[cdrex!17] (9.78,-6.04) rectangle (10.6,-6.86);
\node[inner sep=0pt,font=\sffamily\scriptsize,text=cink] at (10.19,-6.45) {0.19};
\fill[cgone!7] (10.64,-6.04) rectangle (11.46,-6.86);
\node[inner sep=0pt,font=\sffamily\scriptsize,text=cink] at (11.05,-6.45) {\textminus 0.09};
\fill[cdrex!17] (11.5,-6.04) rectangle (12.32,-6.86);
\node[inner sep=0pt,font=\sffamily\scriptsize,text=cink] at (11.91,-6.45) {0.20};
\fill[cgone!2] (12.36,-6.04) rectangle (13.18,-6.86);
\node[inner sep=0pt,font=\sffamily\scriptsize,text=cink] at (12.77,-6.45) {\textminus 0.02};
\fill[cgone!7] (13.22,-6.04) rectangle (14.04,-6.86);
\node[inner sep=0pt,font=\sffamily\scriptsize,text=cink] at (13.63,-6.45) {\textminus 0.09};
\fill[cgone!6] (14.08,-6.04) rectangle (14.9,-6.86);
\node[inner sep=0pt,font=\sffamily\scriptsize,text=cink] at (14.49,-6.45) {\textminus 0.07};
\fill[cdrex!50] (14.94,-6.04) rectangle (15.76,-6.86);
\draw[cink,line width=0.7pt] (14.94,-6.04) rectangle (15.76,-6.86);
\node[inner sep=0pt,font=\sffamily\scriptsize,text=cink] at (15.35,-6.45) {0.59};
\node[mut,anchor=north] at (12.34,-7.1) {People's composite of the dimension};
\node[mut,rotate=90,anchor=south] at (0.2,-3.44) {The model's dimension score};
\end{tikzpicture}%

\caption{Correlation of each model's dimension score (rows) with people's composite of each dimension
(columns), over the 160 excerpts, as observed. The framed diagonal is the dimension's own composite. Grey:
negative.}\label{fig:s2cross}
\end{figure}

\paragraph{Its own dimension, not another.} A dimension score should follow people's composite of its own
dimension and not of another. It does, for every dimension and both models (Figure~\ref{fig:s2cross};
Table~\ref{tab:s2cross}): the mean correlation on the diagonal is 0.68 for \jev{} and 0.66 for \drex{}, against
0.16 and 0.15 off it in absolute value. With either model the margin is clear for six dimensions. It is not for
F4, where \jev{}'s own correlation of 0.51 is nearly matched by a correlation of $-0.45$ between its F6 score and
people's F4 composite.

\paragraph{Relations among the dimensions.} Because the two markers of a dimension are rated by different
people, the part of a relation that comes from one rater answering both questions can be taken out
(Table~\ref{tab:s2halo}). Two relations then differ from the models'. In people's reading, direct and organized
argument is strongly opposed to energetic and angry wording ($-0.68$, against $-0.33$ and $-0.29$ in the models'
scores). And care and regard for others is opposed to energetic and angry wording ($-0.53$), where the models'
scores for the two dimensions are unrelated. A rater's overall impression of an excerpt does not colour all of
that rater's answers: what is left of two answers to different questions correlates at 0.04 and 0.01 on average.

\paragraph{Two markers, one thing?} With rater noise divided out, the two markers of F1, F2, F4 and F6 go
together in people's ratings about as much as in the models' or more (0.70 to 0.94; Table~\ref{tab:s2markercorr}).
On F5 they do not: narrative and rebuttal correlate at 0.12 as observed, so that composite averages two nearly
unrelated readings. On F8 the models' scores on hedging and on knowledge limits correlate at 0.79 and 0.76,
while people's means correlate at 0.10: the models see one thing where people do not.

\subsection{A model against an unlimited panel}\label{sec:panel}

\begin{figure}[!htb]\centering
\begin{tikzpicture}[fig]
\mkdot{cpeople}{4.7}{1.13}
\node[leg,anchor=base west] at (4.93,1.04) {one rating of each marker};
\mkdot{cjev}{8.75}{1.13}
\node[leg,anchor=base west] at (8.98,1.04) {jev-1.14.0 (upper)};
\mkdia{cdrex}{11.876}{1.13}
\node[leg,anchor=base west] at (12.106,1.04) {drex-max (lower)};
\draw[grid] (0,0) -- (15.8,0);
\draw[grid] (0,-0.5) -- (15.8,-0.5);
\draw[grid] (0,-1) -- (15.8,-1);
\draw[grid] (0,-1.5) -- (15.8,-1.5);
\draw[grid] (0,-2) -- (15.8,-2);
\draw[grid] (0,-2.5) -- (15.8,-2.5);
\draw[grid] (0,-3) -- (15.8,-3);
\draw[grid] (0,-3.5) -- (15.8,-3.5);
\draw[grid] (0,-4) -- (15.8,-4);
\node[ttl,anchor=base west] at (4.6,0.42) {Correlation with the panel's consensus};
\draw[grid] (4.6,0) -- (4.6,-4);
\node[tick,anchor=north] at (4.6,-4.12) {0};
\draw[grid] (5.88,0) -- (5.88,-4);
\node[tick,anchor=north] at (5.88,-4.12) {0.2};
\draw[grid] (7.16,0) -- (7.16,-4);
\node[tick,anchor=north] at (7.16,-4.12) {0.4};
\draw[grid] (8.44,0) -- (8.44,-4);
\node[tick,anchor=north] at (8.44,-4.12) {0.6};
\draw[grid] (9.72,0) -- (9.72,-4);
\node[tick,anchor=north] at (9.72,-4.12) {0.8};
\draw[grid] (11,0) -- (11,-4);
\node[tick,anchor=north] at (11,-4.12) {1.0};
\draw[axis] (4.6,-4) -- (11,-4);
\node[mut,anchor=north] at (7.8,-4.58) {Correlation with the panel consensus of the dimension};
\node[ttl,anchor=base east] at (13.9,0.42) {Worth, in raters:};
\node[mut,anchor=base east] at (13.9,0.06) {jev-1.14.0};
\node[ttl,anchor=base east] at (15.5,0.42) {};
\node[mut,anchor=base east] at (15.5,0.06) {drex-max};
\node[lbl,anchor=base east] at (4.4,-0.34) {F1 Direct, organized argument};
\mkdot{cpeople}{8.358}{-0.25}
\mkiv{cjev}{9.539}{10.383}{-0.145}
\mkdot{cjev}{9.929}{-0.145}
\mkiv{cdrex}{9.797}{10.584}{-0.355}
\mkdia{cdrex}{10.152}{-0.355}
\node[lbl,anchor=base east] at (13.9,-0.34) {4.3};
\node[lbl,anchor=base east] at (15.5,-0.34) {5.8};
\node[lbl,anchor=base east] at (4.4,-0.84) {F2 Positive tone vs disappointment};
\mkdot{cpeople}{9.262}{-0.75}
\mkiv{cjev}{9.674}{10.276}{-0.645}
\mkdot{cjev}{9.961}{-0.645}
\mkiv{cdrex}{9.392}{10.04}{-0.855}
\mkdia{cdrex}{9.699}{-0.855}
\node[lbl,anchor=base east] at (13.9,-0.84) {2.1};
\node[lbl,anchor=base east] at (15.5,-0.84) {1.5};
\node[lbl,anchor=base east] at (4.4,-1.34) {F3 Self-disclosure, self-criticism};
\mkdot{cpeople}{9.118}{-1.25}
\mkiv{cjev}{9.832}{10.499}{-1.145}
\mkdot{cjev}{10.175}{-1.145}
\mkiv{cdrex}{9.963}{10.46}{-1.355}
\mkdia{cdrex}{10.209}{-1.355}
\node[lbl,anchor=base east] at (13.9,-1.34) {3.2};
\node[lbl,anchor=base east] at (15.5,-1.34) {3.3};
\node[lbl,anchor=base east] at (4.4,-1.84) {F4 Care and regard for others};
\mkdot{cpeople}{8.025}{-1.75}
\mkiv{cjev}{7.329}{9.597}{-1.645}
\mkdot{cjev}{8.624}{-1.645}
\mkiv{cdrex}{7.264}{9.507}{-1.855}
\mkdia{cdrex}{8.59}{-1.855}
\node[lbl,anchor=base east] at (13.9,-1.84) {1.6};
\node[lbl,anchor=base east] at (15.5,-1.84) {1.6};
\node[lbl,anchor=base east] at (4.4,-2.34) {F5 Narrative vs rebuttal};
\mkdot{cpeople}{7.787}{-2.25}
\mkiv{cjev}{9.927}{10.88}{-2.145}
\mkdot{cjev}{10.431}{-2.145}
\mkiv{cdrex}{9.597}{10.602}{-2.355}
\mkdia{cdrex}{10.13}{-2.355}
\node[lbl,anchor=base east] at (13.9,-2.34) {15};
\node[lbl,anchor=base east] at (15.5,-2.34) {8.9};
\node[lbl,anchor=base east] at (4.4,-2.84) {F6 Energetic and angry vs calm};
\mkdot{cpeople}{8.618}{-2.75}
\mkiv{cjev}{10.281}{10.812}{-2.645}
\mkdot{cjev}{10.547}{-2.645}
\mkiv{cdrex}{9.56}{10.262}{-2.855}
\mkdia{cdrex}{9.929}{-2.855}
\node[lbl,anchor=base east] at (13.9,-2.84) {9.7};
\node[lbl,anchor=base east] at (15.5,-2.84) {3.5};
\node[lbl,anchor=base east] at (4.4,-3.34) {F7 Humor, imaginative language};
\mkdot{cpeople}{7.6}{-3.25}
\mkiv{cjev}{8.472}{10.177}{-3.145}
\mkdot{cjev}{9.301}{-3.145}
\mkiv{cdrex}{7.801}{9.63}{-3.355}
\mkdia{cdrex}{8.749}{-3.355}
\node[lbl,anchor=base east] at (13.9,-3.34) {4.2};
\node[lbl,anchor=base east] at (15.5,-3.34) {2.6};
\node[lbl,anchor=base east] at (4.4,-3.84) {F8 Hedging vs certainty};
\mkdot{cpeople}{6.618}{-3.75}
\mkiv{cjev}{9.047}{10.656}{-3.645}
\mkdot{cjev}{10.132}{-3.645}
\mkiv{cdrex}{9.84}{11}{-3.855}
\mkdia{cdrex}{10.917}{-3.855}
\node[lbl,anchor=base east] at (13.9,-3.84) {27};
\node[lbl,anchor=base east] at (15.5,-3.84) {$>$50};
\end{tikzpicture}%

\caption{Correlation with the consensus of an unlimited panel on the dimension: one rating of each marker, and
each model's dimension score with its 95\% interval. Right: the number of raters per question whose mean follows
the panel as closely as the model does.}\label{fig:s2panel}
\end{figure}

Against what an unlimited panel would say about a dimension, the dimension scores correlate at 0.83 (\jev{}) and
0.81 (\drex{}) on average; one rating of each marker correlates at 0.56 (Figure~\ref{fig:s2panel}). The first
study found 0.86 and 0.51 question by question. A model's dimension score is worth roughly two to six raters per
question on F1 to F4 and F7, and more where raters are noisy and the model is not. The composite would reach a
reliability of 0.60 with two raters per question on F2 and F3, three on F6, four on F1 and F4, six on F5 and F7,
and fifteen on F8 (Table~\ref{tab:s2panel}).

\subsection{Levels}\label{sec:levels}

\begin{figure}[!htb]\centering
\begin{tikzpicture}[fig]
\mkdot{cjev}{3}{1.66}
\node[leg,anchor=base west] at (3.23,1.57) {jev-1.14.0 (upper mark)};
\mkdia{cdrex}{6.786}{1.66}
\node[leg,anchor=base west] at (7.016,1.57) {drex-max (lower mark)};
\mkiv{cjev}{10.158}{10.458}{1.66}
\node[leg,anchor=base west] at (10.538,1.57) {95\% interval};
\draw[grid] (0,0) -- (15.2,0);
\draw[grid] (0,-0.5) -- (15.2,-0.5);
\draw[grid] (0,-1) -- (15.2,-1);
\draw[grid] (0,-1.5) -- (15.2,-1.5);
\draw[grid] (0,-2) -- (15.2,-2);
\draw[grid] (0,-2.5) -- (15.2,-2.5);
\draw[grid] (0,-3) -- (15.2,-3);
\draw[grid] (0,-3.5) -- (15.2,-3.5);
\draw[grid] (0,-4) -- (15.2,-4);
\draw[grid] (0,-4.5) -- (15.2,-4.5);
\draw[grid] (0,-5) -- (15.2,-5);
\draw[grid] (0,-5.5) -- (15.2,-5.5);
\draw[grid] (0,-6) -- (15.2,-6);
\draw[grid] (0,-6.5) -- (15.2,-6.5);
\draw[grid] (0,-7) -- (15.2,-7);
\draw[grid] (0,-7.5) -- (15.2,-7.5);
\draw[grid] (0,-8) -- (15.2,-8);
\node[ttl,anchor=base west] at (2.9,0.95) {Model's mean minus people's mean};
\draw[grid] (2.9,0) -- (2.9,-8);
\node[tick,anchor=north] at (2.9,-8.12) {\textminus 1};
\draw[grid] (4.358,0) -- (4.358,-8);
\node[tick,anchor=north] at (4.358,-8.12) {\textminus 0.5};
\draw[grid] (5.817,0) -- (5.817,-8);
\node[tick,anchor=north] at (5.817,-8.12) {0};
\draw[grid] (7.275,0) -- (7.275,-8);
\node[tick,anchor=north] at (7.275,-8.12) {+0.5};
\draw[grid] (8.733,0) -- (8.733,-8);
\node[tick,anchor=north] at (8.733,-8.12) {+1};
\draw[axis] (2.9,-8) -- (9.9,-8);
\node[mut,anchor=north] at (6.4,-8.58) {Points of the 0--4 scale (right of the line: the model rates higher)};
\draw[ref] (5.817,0.12) -- (5.817,-8);
\node[mut,anchor=base west] at (5.877,0.2) {same level as people};
\node[ttl,anchor=base east] at (11.25,0.95) {People};
\node[mut,anchor=base east] at (11.25,0.59) {mean};
\node[ttl,anchor=base east] at (12.8,0.95) {Slope};
\node[mut,anchor=base east] at (12.8,0.59) {jev-1.14.0};
\node[ttl,anchor=base east] at (14.6,0.95) {Slope};
\node[mut,anchor=base east] at (14.6,0.59) {drex-max};
\node[lbl,anchor=base east] at (2.7,-0.34) {Persuasion};
\mkiv{cjev}{5.394}{6.061}{-0.145}
\mkdot{cjev}{5.709}{-0.145}
\mkiv{cdrex}{4.03}{4.561}{-0.355}
\mkdia{cdrex}{4.274}{-0.355}
\node[lbl,anchor=base east] at (11.25,-0.34) {1.28};
\node[lbl,anchor=base east] at (12.8,-0.34) {0.44};
\node[lbl,anchor=base east] at (14.6,-0.34) {0.65};
\node[lbl,anchor=base east] at (2.7,-0.84) {Persistence};
\mkiv{cjev}{6.871}{7.445}{-0.645}
\mkdot{cjev}{7.141}{-0.645}
\mkiv{cdrex}{4.597}{5.102}{-0.855}
\mkdia{cdrex}{4.851}{-0.855}
\node[lbl,anchor=base east] at (11.25,-0.84) {2.13};
\node[lbl,anchor=base east] at (12.8,-0.84) {0.65};
\node[lbl,anchor=base east] at (14.6,-0.84) {0.76};
\node[lbl,anchor=base east] at (2.7,-1.34) {Positive tone};
\mkiv{cjev}{5.459}{6.036}{-1.145}
\mkdot{cjev}{5.763}{-1.145}
\mkiv{cdrex}{6.343}{6.82}{-1.355}
\mkdia{cdrex}{6.573}{-1.355}
\node[lbl,anchor=base east] at (11.25,-1.34) {1.42};
\node[lbl,anchor=base east] at (12.8,-1.34) {0.67};
\node[lbl,anchor=base east] at (14.6,-1.34) {0.96};
\node[lbl,anchor=base east] at (2.7,-1.84) {Disappointment};
\mkiv{cjev}{6.085}{6.703}{-1.645}
\mkdot{cjev}{6.396}{-1.645}
\mkiv{cdrex}{3.998}{4.611}{-1.855}
\mkdia{cdrex}{4.315}{-1.855}
\node[lbl,anchor=base east] at (11.25,-1.84) {1.83};
\node[lbl,anchor=base east] at (12.8,-1.84) {0.63};
\node[lbl,anchor=base east] at (14.6,-1.84) {0.76};
\node[lbl,anchor=base east] at (2.7,-2.34) {Self-disclosure};
\mkiv{cjev}{8.458}{8.994}{-2.145}
\mkdot{cjev}{8.739}{-2.145}
\mkiv{cdrex}{5.654}{6.022}{-2.355}
\mkdia{cdrex}{5.846}{-2.355}
\node[lbl,anchor=base east] at (11.25,-2.34) {1.64};
\node[lbl,anchor=base east] at (12.8,-2.34) {0.73};
\node[lbl,anchor=base east] at (14.6,-2.34) {0.90};
\node[lbl,anchor=base east] at (2.7,-2.84) {Self-critical};
\mkiv{cjev}{6.135}{6.618}{-2.645}
\mkdot{cjev}{6.368}{-2.645}
\mkiv{cdrex}{4.608}{5.009}{-2.855}
\mkdia{cdrex}{4.808}{-2.855}
\node[lbl,anchor=base east] at (11.25,-2.84) {0.79};
\node[lbl,anchor=base east] at (12.8,-2.84) {0.60};
\node[lbl,anchor=base east] at (14.6,-2.84) {0.81};
\node[lbl,anchor=base east] at (2.7,-3.34) {Compassion};
\mkiv{cjev}{4.792}{5.211}{-3.145}
\mkdot{cjev}{5.009}{-3.145}
\mkiv{cdrex}{4.467}{4.862}{-3.355}
\mkdia{cdrex}{4.674}{-3.355}
\node[lbl,anchor=base east] at (11.25,-3.34) {0.65};
\node[lbl,anchor=base east] at (12.8,-3.34) {0.67};
\node[lbl,anchor=base east] at (14.6,-3.34) {1.14};
\node[lbl,anchor=base east] at (2.7,-3.84) {Respect\textsuperscript{n}};
\mkiv{cjev}{5.434}{5.869}{-3.645}
\mkdot{cjev}{5.651}{-3.645}
\mkiv{cdrex}{5.108}{5.496}{-3.855}
\mkdia{cdrex}{5.309}{-3.855}
\node[lbl,anchor=base east] at (11.25,-3.84) {0.61};
\node[lbl,anchor=base east] at (12.8,-3.84) {0.38};
\node[lbl,anchor=base east] at (14.6,-3.84) {0.58};
\node[lbl,anchor=base east] at (2.7,-4.34) {Narrative};
\mkiv{cjev}{6.568}{7.154}{-4.145}
\mkdot{cjev}{6.869}{-4.145}
\mkiv{cdrex}{4.873}{5.453}{-4.355}
\mkdia{cdrex}{5.165}{-4.355}
\node[lbl,anchor=base east] at (11.25,-4.34) {1.72};
\node[lbl,anchor=base east] at (12.8,-4.34) {0.48};
\node[lbl,anchor=base east] at (14.6,-4.34) {0.52};
\node[lbl,anchor=base east] at (2.7,-4.84) {Rebuttal\textsuperscript{n}};
\mkiv{cjev}{5.48}{6.101}{-4.645}
\mkdot{cjev}{5.788}{-4.645}
\mkiv{cdrex}{4.568}{5.114}{-4.855}
\mkdia{cdrex}{4.843}{-4.855}
\node[lbl,anchor=base east] at (11.25,-4.84) {1.02};
\node[lbl,anchor=base east] at (12.8,-4.84) {0.42};
\node[lbl,anchor=base east] at (14.6,-4.84) {0.51};
\node[lbl,anchor=base east] at (2.7,-5.34) {Energetic};
\mkiv{cjev}{7.31}{7.735}{-5.145}
\mkdot{cjev}{7.524}{-5.145}
\mkiv{cdrex}{4.077}{4.484}{-5.355}
\mkdia{cdrex}{4.276}{-5.355}
\node[lbl,anchor=base east] at (11.25,-5.34) {1.92};
\node[lbl,anchor=base east] at (12.8,-5.34) {0.69};
\node[lbl,anchor=base east] at (14.6,-5.34) {0.95};
\node[lbl,anchor=base east] at (2.7,-5.84) {Calm};
\mkiv{cjev}{4.641}{5.07}{-5.645}
\mkdot{cjev}{4.868}{-5.645}
\mkiv{cdrex}{4.49}{4.901}{-5.855}
\mkdia{cdrex}{4.695}{-5.855}
\node[lbl,anchor=base east] at (11.25,-5.84) {1.23};
\node[lbl,anchor=base east] at (12.8,-5.84) {0.77};
\node[lbl,anchor=base east] at (14.6,-5.84) {1.44};
\node[lbl,anchor=base east] at (2.7,-6.34) {Humor};
\mkiv{cjev}{7.167}{7.643}{-6.145}
\mkdot{cjev}{7.4}{-6.145}
\mkiv{cdrex}{6.304}{6.706}{-6.355}
\mkdia{cdrex}{6.498}{-6.355}
\node[lbl,anchor=base east] at (11.25,-6.34) {0.62};
\node[lbl,anchor=base east] at (12.8,-6.34) {0.43};
\node[lbl,anchor=base east] at (14.6,-6.34) {0.51};
\node[lbl,anchor=base east] at (2.7,-6.84) {Imaginative\textsuperscript{n}};
\mkiv{cjev}{4.509}{5.103}{-6.645}
\mkdot{cjev}{4.811}{-6.645}
\mkiv{cdrex}{4.025}{4.47}{-6.855}
\mkdia{cdrex}{4.26}{-6.855}
\node[lbl,anchor=base east] at (11.25,-6.84) {1.22};
\node[lbl,anchor=base east] at (12.8,-6.84) {0.29};
\node[lbl,anchor=base east] at (14.6,-6.84) {0.47};
\node[lbl,anchor=base east] at (2.7,-7.34) {Hedging};
\mkiv{cjev}{8.662}{9.166}{-7.145}
\mkdot{cjev}{8.91}{-7.145}
\mkiv{cdrex}{5.839}{6.282}{-7.355}
\mkdia{cdrex}{6.044}{-7.355}
\node[lbl,anchor=base east] at (11.25,-7.34) {1.35};
\node[lbl,anchor=base east] at (12.8,-7.34) {0.47};
\node[lbl,anchor=base east] at (14.6,-7.34) {0.61};
\node[lbl,anchor=base east] at (2.7,-7.84) {Knowledge limits\textsuperscript{n}};
\mkiv{cjev}{6.625}{7.394}{-7.645}
\mkdot{cjev}{7.023}{-7.645}
\mkiv{cdrex}{4.108}{4.712}{-7.855}
\mkdia{cdrex}{4.407}{-7.855}
\node[lbl,anchor=base east] at (11.25,-7.84) {1.40};
\node[lbl,anchor=base east] at (12.8,-7.84) {0.07};
\node[lbl,anchor=base east] at (14.6,-7.84) {0.16};
\end{tikzpicture}%

\caption{Each model's mean score minus people's mean rating over the 160 excerpts of the second study, with 95\%
intervals, and how many scale points people's mean rating rises per point of the model's score.
\textsuperscript{n}: not rated in the first study.}\label{fig:s2levels}
\end{figure}

The models' levels are not people's (Figure~\ref{fig:s2levels}). \jev{} rates hedging 1.06 scale points higher
than people and self-disclosure 1.00 higher; over the sixteen questions it is 0.37 points away on average.
\drex{} rates twelve of the sixteen questions lower than people. Both findings repeat the first study. A
model's score cannot be read as a level on people's scale.

\subsection{Videos and streams, earlier and later years}\label{sec:groups}

\begin{figure}[!htb]\centering
\begin{tikzpicture}[fig]
\mkdot{cjev}{1}{7.39}
\node[leg,anchor=base west] at (1.23,7.3) {jev-1.14.0};
\mkdia{cdrex}{3.07}{7.39}
\node[leg,anchor=base west] at (3.3,7.3) {drex-max};
\mkiv{cmuted}{4.726}{5.026}{7.39}
\node[leg,anchor=base west] at (5.106,7.3) {the same question};
\draw[grid] (0.9,0) -- (0.9,6);
\node[tick,anchor=north] at (0.9,-0.1) {\textminus 1.5};
\draw[grid] (1.9,0) -- (1.9,6);
\node[tick,anchor=north] at (1.9,-0.1) {\textminus 1.0};
\draw[grid] (2.9,0) -- (2.9,6);
\node[tick,anchor=north] at (2.9,-0.1) {\textminus 0.5};
\draw[grid] (3.9,0) -- (3.9,6);
\node[tick,anchor=north] at (3.9,-0.1) {0};
\draw[grid] (4.9,0) -- (4.9,6);
\node[tick,anchor=north] at (4.9,-0.1) {+0.5};
\draw[grid] (5.9,0) -- (5.9,6);
\node[tick,anchor=north] at (5.9,-0.1) {+1.0};
\draw[grid] (6.9,0) -- (6.9,6);
\node[tick,anchor=north] at (6.9,-0.1) {+1.5};
\draw[grid] (0.9,0) -- (6.9,0);
\node[tick,anchor=east] at (0.8,0) {\textminus 1.5};
\draw[grid] (0.9,1) -- (6.9,1);
\node[tick,anchor=east] at (0.8,1) {\textminus 1.0};
\draw[grid] (0.9,2) -- (6.9,2);
\node[tick,anchor=east] at (0.8,2) {\textminus 0.5};
\draw[grid] (0.9,3) -- (6.9,3);
\node[tick,anchor=east] at (0.8,3) {0};
\draw[grid] (0.9,4) -- (6.9,4);
\node[tick,anchor=east] at (0.8,4) {+0.5};
\draw[grid] (0.9,5) -- (6.9,5);
\node[tick,anchor=east] at (0.8,5) {+1.0};
\draw[grid] (0.9,6) -- (6.9,6);
\node[tick,anchor=east] at (0.8,6) {+1.5};
\draw[axis] (0.9,0) -- (6.9,0);
\draw[axis] (0.9,0) -- (0.9,6);
\node[mut,anchor=north] at (3.9,-0.55) {Model: streams minus videos};
\node[mut,rotate=90,anchor=south] at (-0.08,3) {People: streams minus videos};
\node[ttl,anchor=base west] at (0.9,6.48) {Streams minus videos};
\node[mut,anchor=base west] at (0.9,6.13) {$r$ with jev-1.14.0: 0.86; with drex-max: 0.84};
\draw[axis] (3.9,0) -- (3.9,6);
\draw[axis] (0.9,3) -- (6.9,3);
\draw[ref] (0.9,0) -- (6.9,6);
\node[mut,rotate=45,anchor=south east] at (6.75,5.91) {people = model};
\draw[cmuted,line width=0.5pt] (2.833,2.301) -- (3.279,2.301);
\mkdot{cjev}{2.833}{2.301}
\mkdia{cdrex}{3.279}{2.301}
\draw[cmuted,line width=0.5pt] (2.762,1.811) -- (3.235,1.811);
\mkdot{cjev}{2.762}{1.811}
\mkdia{cdrex}{3.235}{1.811}
\draw[cmuted,line width=0.5pt] (2.973,2.147) -- (3.559,2.147);
\mkdot{cjev}{2.973}{2.147}
\mkdia{cdrex}{3.559}{2.147}
\draw[cmuted,line width=0.5pt] (4.514,3.617) -- (4.05,3.617);
\mkdot{cjev}{4.514}{3.617}
\mkdia{cdrex}{4.05}{3.617}
\draw[cmuted,line width=0.5pt] (4.364,3.46) -- (4.364,3.46);
\mkdot{cjev}{4.364}{3.46}
\mkdia{cdrex}{4.364}{3.46}
\draw[cmuted,line width=0.5pt] (5.018,3.623) -- (4.361,3.623);
\mkdot{cjev}{5.018}{3.623}
\mkdia{cdrex}{4.361}{3.623}
\draw[cmuted,line width=0.5pt] (4.162,2.823) -- (3.928,2.823);
\mkdot{cjev}{4.162}{2.823}
\mkdia{cdrex}{3.928}{2.823}
\draw[cmuted,line width=0.5pt] (4.056,2.669) -- (3.984,2.669);
\mkdot{cjev}{4.056}{2.669}
\mkdia{cdrex}{3.984}{2.669}
\draw[cmuted,line width=0.5pt] (3.638,1.88) -- (3.524,1.88);
\mkdot{cjev}{3.638}{1.88}
\mkdia{cdrex}{3.524}{1.88}
\draw[cmuted,line width=0.5pt] (4.263,3.123) -- (4.3,3.123);
\mkdot{cjev}{4.263}{3.123}
\mkdia{cdrex}{4.3}{3.123}
\draw[cmuted,line width=0.5pt] (4.524,3.625) -- (3.972,3.625);
\mkdot{cjev}{4.524}{3.625}
\mkdia{cdrex}{3.972}{3.625}
\draw[cmuted,line width=0.5pt] (2.817,1.423) -- (3.244,1.423);
\mkdot{cjev}{2.817}{1.423}
\mkdia{cdrex}{3.244}{1.423}
\draw[cmuted,line width=0.5pt] (4.885,3.498) -- (4.766,3.498);
\mkdot{cjev}{4.885}{3.498}
\mkdia{cdrex}{4.766}{3.498}
\draw[cmuted,line width=0.5pt] (4.447,2.768) -- (4.109,2.768);
\mkdot{cjev}{4.447}{2.768}
\mkdia{cdrex}{4.109}{2.768}
\draw[cmuted,line width=0.5pt] (4.376,3.474) -- (4.137,3.474);
\mkdot{cjev}{4.376}{3.474}
\mkdia{cdrex}{4.137}{3.474}
\draw[cmuted,line width=0.5pt] (4.419,2.471) -- (4.102,2.471);
\mkdot{cjev}{4.419}{2.471}
\mkdia{cdrex}{4.102}{2.471}
\node[leg,anchor=base west] at (3.384,1.343) {Calm};
\node[leg,anchor=base west] at (3.778,1.8) {Narrative};
\node[leg,anchor=base east] at (2.693,2.221) {Persuasion};
\node[leg,anchor=base east] at (3.832,3.545) {Energetic};
\draw[grid] (9,0) -- (9,6);
\node[tick,anchor=north] at (9,-0.1) {\textminus 0.2};
\draw[grid] (10.5,0) -- (10.5,6);
\node[tick,anchor=north] at (10.5,-0.1) {\textminus 0.1};
\draw[grid] (12,0) -- (12,6);
\node[tick,anchor=north] at (12,-0.1) {0};
\draw[grid] (13.5,0) -- (13.5,6);
\node[tick,anchor=north] at (13.5,-0.1) {+0.1};
\draw[grid] (15,0) -- (15,6);
\node[tick,anchor=north] at (15,-0.1) {+0.2};
\draw[grid] (9,0) -- (15,0);
\node[tick,anchor=east] at (8.9,0) {\textminus 0.2};
\draw[grid] (9,1.5) -- (15,1.5);
\node[tick,anchor=east] at (8.9,1.5) {\textminus 0.1};
\draw[grid] (9,3) -- (15,3);
\node[tick,anchor=east] at (8.9,3) {0};
\draw[grid] (9,4.5) -- (15,4.5);
\node[tick,anchor=east] at (8.9,4.5) {+0.1};
\draw[grid] (9,6) -- (15,6);
\node[tick,anchor=east] at (8.9,6) {+0.2};
\draw[axis] (9,0) -- (15,0);
\draw[axis] (9,0) -- (9,6);
\node[mut,anchor=north] at (12,-0.55) {Model: change per year};
\node[mut,rotate=90,anchor=south] at (8.02,3) {People: change per year};
\node[ttl,anchor=base west] at (9,6.48) {Change per year};
\node[mut,anchor=base west] at (9,6.13) {$r$ with jev-1.14.0: 0.87; with drex-max: 0.87};
\draw[axis] (12,0) -- (12,6);
\draw[axis] (9,3) -- (15,3);
\draw[ref] (9,0) -- (15,6);
\node[mut,rotate=45,anchor=south east] at (14.85,5.91) {people = model};
\draw[cmuted,line width=0.5pt] (12.646,3.348) -- (12.305,3.348);
\mkdot{cjev}{12.646}{3.348}
\mkdia{cdrex}{12.305}{3.348}
\draw[cmuted,line width=0.5pt] (11.413,2.403) -- (11.587,2.403);
\mkdot{cjev}{11.413}{2.403}
\mkdia{cdrex}{11.587}{2.403}
\draw[cmuted,line width=0.5pt] (13.466,3.74) -- (12.713,3.74);
\mkdot{cjev}{13.466}{3.74}
\mkdia{cdrex}{12.713}{3.74}
\draw[cmuted,line width=0.5pt] (12.256,2.218) -- (12.09,2.218);
\mkdot{cjev}{12.256}{2.218}
\mkdia{cdrex}{12.09}{2.218}
\draw[cmuted,line width=0.5pt] (9.797,0.954) -- (9.787,0.954);
\mkdot{cjev}{9.797}{0.954}
\mkdia{cdrex}{9.787}{0.954}
\draw[cmuted,line width=0.5pt] (10.81,1.943) -- (11.25,1.943);
\mkdot{cjev}{10.81}{1.943}
\mkdia{cdrex}{11.25}{1.943}
\draw[cmuted,line width=0.5pt] (11.391,2.12) -- (11.582,2.12);
\mkdot{cjev}{11.391}{2.12}
\mkdia{cdrex}{11.582}{2.12}
\draw[cmuted,line width=0.5pt] (11.555,3.043) -- (11.594,3.043);
\mkdot{cjev}{11.555}{3.043}
\mkdia{cdrex}{11.594}{3.043}
\draw[cmuted,line width=0.5pt] (11.93,2.481) -- (11.409,2.481);
\mkdot{cjev}{11.93}{2.481}
\mkdia{cdrex}{11.409}{2.481}
\draw[cmuted,line width=0.5pt] (11.442,2.745) -- (11.56,2.745);
\mkdot{cjev}{11.442}{2.745}
\mkdia{cdrex}{11.56}{2.745}
\draw[cmuted,line width=0.5pt] (12.604,3.524) -- (12.667,3.524);
\mkdot{cjev}{12.604}{3.524}
\mkdia{cdrex}{12.667}{3.524}
\draw[cmuted,line width=0.5pt] (11.895,2.802) -- (11.715,2.802);
\mkdot{cjev}{11.895}{2.802}
\mkdia{cdrex}{11.715}{2.802}
\draw[cmuted,line width=0.5pt] (12.021,2.667) -- (11.821,2.667);
\mkdot{cjev}{12.021}{2.667}
\mkdia{cdrex}{11.821}{2.667}
\draw[cmuted,line width=0.5pt] (11.46,2.721) -- (11.544,2.721);
\mkdot{cjev}{11.46}{2.721}
\mkdia{cdrex}{11.544}{2.721}
\draw[cmuted,line width=0.5pt] (11.799,2.755) -- (11.492,2.755);
\mkdot{cjev}{11.799}{2.755}
\mkdia{cdrex}{11.492}{2.755}
\draw[cmuted,line width=0.5pt] (11.886,3.103) -- (11.624,3.103);
\mkdot{cjev}{11.886}{3.103}
\mkdia{cdrex}{11.624}{3.103}
\node[leg,anchor=base west] at (9.937,0.874) {Self-disclosure};
\node[leg,anchor=base west] at (12.396,2.138) {Disappointment};
\node[leg,anchor=base east] at (11.273,2.323) {Persistence};
\node[leg,anchor=base east] at (12.165,3.268) {Persuasion};
\end{tikzpicture}%

\caption{Differences between groups of excerpts in people's mean ratings (vertical) and in each model's scores
(horizontal), one pair of marks per question. Left: streams minus videos. Right: change per year, from a
regression on the kind and the year of the recording.}\label{fig:s2groups}
\end{figure}

\paragraph{People see the differences the models see.} Which questions are rated higher in streams and which in
videos is nearly the same for people and models: across the twelve questions of the first study the differences
correlate at 0.92 with \jev{} in each study, and across all sixteen at 0.86 (Figure~\ref{fig:s2groups};
Table~\ref{tab:s2groupsacross}). The changes over the years agree as well. Self-disclosure falls by 0.14 scale
points a year in people's ratings and by 0.15 in both models'.

\begin{figure}[!htb]\centering
\begin{tikzpicture}[fig]
\mkdot{cjev}{4.7}{1.66}
\node[leg,anchor=base west] at (4.93,1.57) {jev-1.14.0 (upper mark)};
\mkdia{cdrex}{8.486}{1.66}
\node[leg,anchor=base west] at (8.716,1.57) {drex-max (lower mark)};
\mkiv{cjev}{11.858}{12.158}{1.66}
\node[leg,anchor=base west] at (12.238,1.57) {95\% interval};
\draw[grid] (0,0) -- (14.8,0);
\draw[grid] (0,-0.5) -- (14.8,-0.5);
\draw[grid] (0,-1) -- (14.8,-1);
\draw[grid] (0,-1.5) -- (14.8,-1.5);
\draw[grid] (0,-2) -- (14.8,-2);
\draw[grid] (0,-2.5) -- (14.8,-2.5);
\draw[grid] (0,-3) -- (14.8,-3);
\draw[grid] (0,-3.5) -- (14.8,-3.5);
\draw[grid] (0,-4) -- (14.8,-4);
\node[ttl,anchor=base west] at (4.6,0.95) {Streams against videos};
\draw[grid] (4.6,0) -- (4.6,-4);
\node[tick,anchor=north] at (4.6,-4.12) {\textminus 1.0};
\draw[grid] (5.75,0) -- (5.75,-4);
\node[tick,anchor=north] at (5.75,-4.12) {\textminus 0.5};
\draw[grid] (6.9,0) -- (6.9,-4);
\node[tick,anchor=north] at (6.9,-4.12) {0};
\draw[grid] (8.05,0) -- (8.05,-4);
\node[tick,anchor=north] at (8.05,-4.12) {+0.5};
\draw[grid] (9.2,0) -- (9.2,-4);
\node[tick,anchor=north] at (9.2,-4.12) {+1.0};
\draw[axis] (4.6,-4) -- (9.2,-4);
\node[mut,anchor=north] at (6.9,-4.58) {Standard deviations of people's composite};
\draw[ref] (6.9,0.12) -- (6.9,-4);
\node[ttl,anchor=base west] at (10.1,0.95) {Per year};
\draw[grid] (10.867,0) -- (10.867,-4);
\node[tick,anchor=north] at (10.867,-4.12) {\textminus 0.1};
\draw[grid] (12.4,0) -- (12.4,-4);
\node[tick,anchor=north] at (12.4,-4.12) {0};
\draw[grid] (13.933,0) -- (13.933,-4);
\node[tick,anchor=north] at (13.933,-4.12) {+0.1};
\draw[axis] (10.1,-4) -- (14.7,-4);
\node[mut,anchor=north] at (12.4,-4.58) {Standard deviations per year};
\draw[ref] (12.4,0.12) -- (12.4,-4);
\node[lbl,anchor=base east] at (4.4,-0.34) {F1 Direct, organized argument};
\mkiv{cjev}{5.97}{7.253}{-0.145}
\mkdot{cjev}{6.634}{-0.145}
\mkiv{cdrex}{5.976}{7.023}{-0.355}
\mkdia{cdrex}{6.488}{-0.355}
\mkiv{cjev}{11.566}{12.979}{-0.145}
\mkdot{cjev}{12.275}{-0.145}
\mkiv{cdrex}{11.899}{13.145}{-0.355}
\mkdia{cdrex}{12.526}{-0.355}
\node[lbl,anchor=base east] at (4.4,-0.84) {F2 Positive tone vs disappointment};
\mkiv{cjev}{5.608}{6.675}{-0.645}
\mkdot{cjev}{6.153}{-0.645}
\mkiv{cdrex}{5.555}{6.664}{-0.855}
\mkdia{cdrex}{6.082}{-0.855}
\mkiv{cjev}{11.747}{13.023}{-0.645}
\mkdot{cjev}{12.38}{-0.645}
\mkiv{cdrex}{11.969}{13.282}{-0.855}
\mkdia{cdrex}{12.607}{-0.855}
\node[lbl,anchor=base east] at (4.4,-1.34) {F3 Self-disclosure, self-criticism};
\mkiv{cjev}{6.705}{7.688}{-1.145}
\mkdot{cjev}{7.207}{-1.145}
\mkiv{cdrex}{6.932}{7.884}{-1.355}
\mkdia{cdrex}{7.401}{-1.355}
\mkiv{cjev}{11.11}{12.316}{-1.145}
\mkdot{cjev}{11.7}{-1.145}
\mkiv{cdrex}{11.393}{12.548}{-1.355}
\mkdia{cdrex}{11.973}{-1.355}
\node[lbl,anchor=base east] at (4.4,-1.84) {F4 Care and regard for others};
\mkiv{cjev}{5.005}{6.472}{-1.645}
\mkdot{cjev}{5.744}{-1.645}
\mkiv{cdrex}{5.048}{6.639}{-1.855}
\mkdia{cdrex}{5.86}{-1.855}
\mkiv{cjev}{11.248}{12.814}{-1.645}
\mkdot{cjev}{11.996}{-1.645}
\mkiv{cdrex}{11.546}{13.171}{-1.855}
\mkdia{cdrex}{12.353}{-1.855}
\node[lbl,anchor=base east] at (4.4,-2.34) {F5 Narrative vs rebuttal};
\mkiv{cjev}{6.155}{7.282}{-2.145}
\mkdot{cjev}{6.733}{-2.145}
\mkiv{cdrex}{5.98}{7.19}{-2.355}
\mkdia{cdrex}{6.605}{-2.355}
\mkiv{cjev}{10.927}{12.599}{-2.145}
\mkdot{cjev}{11.822}{-2.145}
\mkiv{cdrex}{11.123}{12.941}{-2.355}
\mkdia{cdrex}{12.05}{-2.355}
\node[lbl,anchor=base east] at (4.4,-2.84) {F6 Energetic and angry vs calm};
\mkiv{cjev}{6.791}{7.938}{-2.645}
\mkdot{cjev}{7.367}{-2.645}
\mkiv{cdrex}{7.171}{8.456}{-2.855}
\mkdia{cdrex}{7.833}{-2.855}
\mkiv{cjev}{12.293}{13.462}{-2.645}
\mkdot{cjev}{12.829}{-2.645}
\mkiv{cdrex}{12.211}{13.543}{-2.855}
\mkdia{cdrex}{12.816}{-2.855}
\node[lbl,anchor=base east] at (4.4,-3.34) {F7 Humor, imaginative language};
\mkiv{cjev}{6.324}{7.548}{-3.145}
\mkdot{cjev}{6.928}{-3.145}
\mkiv{cdrex}{6.489}{7.776}{-3.355}
\mkdia{cdrex}{7.14}{-3.355}
\mkiv{cjev}{10.603}{12.053}{-3.145}
\mkdot{cjev}{11.313}{-3.145}
\mkiv{cdrex}{10.745}{12.317}{-3.355}
\mkdia{cdrex}{11.54}{-3.355}
\node[lbl,anchor=base east] at (4.4,-3.84) {F8 Hedging vs certainty};
\mkiv{cjev}{5.758}{7.189}{-3.645}
\mkdot{cjev}{6.46}{-3.645}
\mkiv{cdrex}{6.094}{7.441}{-3.855}
\mkdia{cdrex}{6.739}{-3.855}
\mkiv{cjev}{11.208}{12.805}{-3.645}
\mkdot{cjev}{12.025}{-3.645}
\mkiv{cdrex}{11.386}{12.984}{-3.855}
\mkdia{cdrex}{12.18}{-3.855}
\end{tikzpicture}%

\caption{At the same model score, how much higher or lower people rate streams than videos (left) and later
years than earlier ones (right), in standard deviations of people's composite; from a regression of the
composite on the model's dimension score, the kind and the year of the recording. 95\% intervals. Zero: the
model's score means the same in both groups.}\label{fig:s2bias}
\end{figure}

\paragraph{Does a score mean the same in both groups?} The site compares videos with streams and earlier with
later years. Such a comparison is fair when a given score means the same in both groups
(Figure~\ref{fig:s2bias}; Table~\ref{tab:s2bias}). On F1, F5 and F8 it does, for both models and both
comparisons. On F4, care and regard for others, it does not: at the same score people rate streams half a
standard deviation lower than videos, $-0.50$ ($-0.82$ to $-0.19$) for \jev{} and $-0.45$ ($-0.81$ to $-0.11$)
for \drex{}, so the models' F4 scores overstate streams against videos. On F2 people rate streams a third of a
standard deviation lower at the same score: the models see streams as less positive than videos, and people see
a larger difference. On F7 people's rating falls by 0.07 standard deviations a year at the same score. With 32
comparisons one or two intervals would leave out zero by chance.

\subsection{Both models together, and the collection}\label{sec:both}

Averaging the two models' dimension scores gives 0.681 over the eight dimensions against 0.670 for \jev{} alone.
On no dimension is the average clearly better than the better single model (Table~\ref{tab:s2both}), because
the two models miss the same excerpts: with rater noise taken out, their misses correlate at 0.53 to 0.94. A
second model adds little because it shares the first one's blind spots.

Two checks of the collection. Over the eight excerpts of a session a rating does not get noisier: the last four
excerpts against the first four give a ratio of 0.96 (0.91 to 1.02). And the question shown first beside the
excerpt is not rated differently from the others: being first shifts a question's level by $+0.01$ scale points
on average (standard error 0.05).
